\kefalaio{Οι φυσικοί αριθμοί}{Από το λεξιλόγιο στη γραμματική}{%
\begin{tabular}{@{}l@{\hspace{1em}}p{0.82\linewidth}@{}}
1.1 & Φυσικοί αριθμοί, διάταξη, στρογγυλοποίηση\\[2pt]
1.2 & Πρόσθεση, αφαίρεση και πολλαπλασιασμός φυσικών αριθμών\\[2pt]
1.3 & Δυνάμεις φυσικών αριθμών\\[2pt]
1.4 & Ευκλείδεια διαίρεση, διαιρετότητα\\[2pt]
1.5 & Χαρακτήρες διαιρετότητας, ΜΚΔ, ΕΚΠ, ανάλυση σε γινόμενο πρώτων παραγόντων\\[2pt]
 & Επανάληψη του κεφαλαίου\\
\end{tabular}}
\paragrafos{Φυσικοί αριθμοί, διάταξη, στρογγυλοποίηση: η σειρά που δεν τελειώνει}{Α.1.1, «Φυσικοί αριθμοί – Διάταξη φυσικών – Στρογγυλοποίηση»}

\begin{anagke}
Πρώτη μέρα στο γυμνάσιο. Στην αίθουσα της Α3 η Ζωή, ο Ηλίας, η Μυρτώ, ο Στέλιος, η Φανή και ο Βασίλης, ο πρώην πρόεδρος του μαθητικού συμβουλίου της ΣΤ’, βρίσκουν στα θρανία δύο γνωστά πρόσωπα από το άλλο δημοτικό της περιοχής: τον Νίκο και την Αγγελική, που τους έχουν δει σε αγώνες μπάσκετ και σε γιορτές. Στο τελευταίο θρανίο κάθεται ένα κορίτσι που δεν το ξέρει κανείς. Την λένε Πολυξένη, ήρθε το καλοκαίρι από τη Θεσσαλονίκη και δεν έχει πει ακόμα ούτε λέξη.

Στην είσοδο ένα χαρτί λέει: «Το σχολείο μας έχει 387 μαθητές». Ο Νίκος το φωτογραφίζει και στέλνει στη μαμά του: «Το σχολείο έχει 400 παιδιά!». Η μαμά του απαντάει με ένα χαμογελαστό «στρογγύλεψες».

Η κυρία Αλεξάνδρα, η καθηγήτρια των μαθηματικών, δεν ανοίγει βιβλίο. Γράφει στον πίνακα δύο αριθμούς: τους μαθητές του σχολείου, 387, και το υψόμετρο του Ολύμπου, 2.917 μέτρα. «Έχετε να τους πείτε σε έναν φίλο σας. Ποιον θα πείτε ακριβώς και ποιον περίπου;»

«Όλους ακριβώς», λέει ο Στέλιος. «Τι νόημα έχει το “περίπου”;» Η Ζωή σηκώνει το χέρι. «Αν πεις στον φίλο σου ότι περπάτησες 8.412 βήματα, θα θυμάται μόνο το 8 και θα μπερδεύεται με τα υπόλοιπα. Το “περίπου 8.000” το κρατάς». Ο Βασίλης χαμογελάει: «Εγώ λέω πάντα “γύρω στα τετρακόσια”. Με το μάτι». Τότε ακούγεται για πρώτη φορά η φωνή της Πολυξένης, χαμηλά: «Και πώς ξέρουμε ότι όλοι θα βγάλουμε το ίδιο “περίπου”; Αν ο Νίκος κάνει το 387 τετρακόσια, και ο Βασίλης “γύρω στα τριακόσια ογδόντα”;»

Όλοι γυρίζουν να την κοιτάξουν. Η κυρία Αλεξάνδρα χαμογελάει: η ερώτηση είναι ακριβώς αυτή που ήθελε να ακούσει. Πότε επιτρέπεται το «περίπου», και πώς γίνεται με κανόνα, ώστε όλοι να κάνουμε το ίδιο;
\end{anagke}

\begin{lexi}{φυσικός, στρογγυλοποίηση}{φυσικός, στρογγυλοποίηση}
Οι \lexiword{φυσικοί} αριθμοί λέγονται έτσι επειδή είναι οι πρώτοι που βρίσκει κανείς στη φύση, όταν μετράει πράγματα: μήλα, παιδιά, βήματα. Κανείς δεν μέτρησε ποτέ «ενάμισι παιδί». Το μηδέν μπήκε στην παρέα αργότερα. Στο σχολείο θα το βρεις μέσα στους φυσικούς.

Η \lexiword{στρογγυλοποίηση} είναι το να κάνεις έναν αριθμό «στρογγυλό», δηλαδή να τον αλλάξεις με έναν κοντινό του που τελειώνει σε μηδενικά: 387 γίνεται 400. Στρογγυλός σημαίνει τακτοποιημένος, χωρίς αιχμές. Δεν είναι μαγικό: είναι \emph{συμφωνία}. Όλοι κάνουμε τα ίδια βήματα και βγάζουμε το ίδιο αποτέλεσμα.

Δύο ακόμα λέξεις που θα ξανασυναντήσεις: \lexiword{άρτιος} είναι ο αριθμός που «ταιριάζει» τέλεια σε ζευγάρια, και \lexiword{περιττός} αυτός που «περισσεύει» ένας. Τους λέμε και ζυγούς και μονούς.
\end{lexi}

\begin{metafrasi}
Ο αριθμός που ακολουθεί το 9. & $9+1=10$\\
Το πέντε είναι μεγαλύτερο από το τρία. & $5>3$\\
Ο αριθμός $x$ βρίσκεται ανάμεσα στο 3 και στο 8. & $3<x<8$\\
Τρία εκατομμύρια, σαράντα χιλιάδες και επτά. & 3.040.007\\
Το 3.486.917 στρογγυλοποιημένο στις χιλιάδες. & $3.486.917\approx 3.487.000$\\
Ο αριθμός είναι άρτιος. & Διαιρείται με το 2\\
\end{metafrasi}

\tmima{Η θεωρία}

\begin{orismos}{φυσικοί αριθμοί}
Οι αριθμοί $0,\,1,\,2,\,3,\,4,\,\dots,\,99,\,100,\,\dots$ λέγονται \textbf{φυσικοί αριθμοί}. Κάθε φυσικός έχει έναν επόμενο και, εκτός από το 0, έναν προηγούμενο.
\end{orismos}

Οι φυσικοί χωρίζονται σε δύο κατηγορίες: \textbf{άρτιους} (ή ζυγούς), που διαιρούνται με το 2, και \textbf{περιττούς} (ή μονούς), που δεν διαιρούνται. Ο άρτιος τελειώνει πάντα σε 0, 2, 4, 6 ή 8.

\tmimasmall{Δέκα ψηφία, άπειροι αριθμοί}

Με τα δέκα ψηφία $0,\,1,\,2,\,\dots,\,9$ γράφουμε όσους αριθμούς θέλουμε. Αυτό γίνεται επειδή η αξία ενός ψηφίου εξαρτάται και από τη \textbf{θέση} του, δηλαδή από τη \textbf{δεκαδική τάξη}: μονάδες, δεκάδες, εκατοντάδες, χιλιάδες, δεκάδες χιλιάδων, εκατοντάδες χιλιάδων, εκατομμύρια. Το 0 κρατάει μια θέση άδεια: στο 6.040.205 τα τρία μηδενικά λένε ότι δεν υπάρχουν εκατοντάδες χιλιάδων, χιλιάδες και δεκάδες.

\tmimasmall{Διάταξη: ποιος είναι μεγαλύτερος}

Γράφουμε $=$ για «ίσος με», $<$ για «μικρότερος από» και $>$ για «μεγαλύτερος από». Δύο φυσικούς μπορούμε πάντα να τους συγκρίνουμε.

\begin{kanonas}{πώς συγκρίνω δύο φυσικούς}
\begin{enumerate}[leftmargin=1.6em,itemsep=0.2ex,topsep=2pt]
\item Μετράω τα ψηφία. Ο αριθμός με περισσότερα ψηφία είναι μεγαλύτερος (αν δεν ξεκινάει με 0).
\item Αν έχουν ίδιο πλήθος ψηφίων, τα κοιτάζω από αριστερά προς τα δεξιά. Στο πρώτο ψηφίο που διαφέρουν, μεγαλύτερος είναι αυτός με το μεγαλύτερο ψηφίο.
\end{enumerate}
\end{kanonas}

Έτσι $73.409>73.094$, γιατί τα δύο πρώτα ψηφία ταιριάζουν και στο τρίτο το 4 είναι μεγαλύτερο από το 0. Και $812<10.250$, γιατί το 812 έχει λιγότερα ψηφία. Τους φυσικούς μπορούμε να τους βάλουμε σε σειρά: \textbf{αύξουσα}, από τον μικρότερο στον μεγαλύτερο, ή \textbf{φθίνουσα}, αντίστροφα. Δεν υπάρχει μεγαλύτερος φυσικός αριθμός, γιατί κάθε φυσικός έχει επόμενο. Ο μικρότερος είναι το 0.

\tmimasmall{Η αριθμητική ευθεία}

Η διάταξη μας επιτρέπει να βάλουμε τους φυσικούς πάνω σε μια ευθεία. Διαλέγουμε ένα σημείο Ο για το 0 (την αρχή) και ένα σημείο Α δεξιά του για το 1. Το τμήμα ΟΑ είναι η \textbf{μονάδα μέτρησης}. Με το ίδιο μήκος βρίσκουμε τα σημεία των 2, 3, 4, \dots

\begin{center}
\begin{tikzpicture}[x=1.15cm,y=1cm]
\draw[-{Stealth}] (-0.3,0)--(9.6,0);
\foreach \i in {0,...,8}{\draw (\i,-0.12)--(\i,0.12); \node[below,font=\small] at (\i,-0.12) {\i};}
\node[above,font=\small\bfseries,gblue] at (0,0.15) {Ο};
\node[above,font=\small\bfseries,gblue] at (1,0.15) {Α};
\draw[<->,gblue] (0,0.75)--(1,0.75) node[midway,above,font=\scriptsize] {μονάδα};
\end{tikzpicture}
\end{center}

\tmimasmall{Στρογγυλοποίηση}

\begin{kanonas}{στρογγυλοποίηση φυσικού}
\begin{enumerate}[leftmargin=1.6em,itemsep=0.2ex,topsep=2pt]
\item Βρίσκω την τάξη στην οποία θα στρογγυλοποιήσω (δεκάδες, εκατοντάδες, χιλιάδες, \dots).
\item Κοιτάζω το ψηφίο της \textbf{αμέσως μικρότερης} τάξης.
\item Αν είναι $0,\,1,\,2,\,3$ ή $4$, το ψηφίο της τάξης μου μένει ίδιο. Αν είναι $5,\,6,\,7,\,8$ ή $9$, αυξάνεται κατά 1.
\item Όλα τα ψηφία δεξιά από την τάξη μου γίνονται 0.
\end{enumerate}
\end{kanonas}

Το αποτέλεσμα λέγεται \textbf{προσέγγιση} του αριθμού και γράφεται με το σύμβολο $\approx$ (διαβάζεται «περίπου ίσο»). Για παράδειγμα, $2.917\approx 3.000$ στις χιλιάδες και $2.917\approx 2.900$ στις εκατοντάδες.

Η στρογγυλοποίηση ταιριάζει εκεί που μας ενδιαφέρει η τάξη μεγέθους, όπως το πλήθος των ανθρώπων σε μια πόλη. Δεν ταιριάζει εκεί που χρειάζεται ακρίβεια: ο αριθμός του τηλεφώνου, το ΑΦΜ, ο κωδικός ενός λουκέτου.

\begin{prosoxi}
\begin{itemize}[leftmargin=1.2em,itemsep=0.3ex,topsep=2pt]
\item Στρογγυλοποιούμε \textbf{μία φορά}, από τον αρχικό αριθμό. Το 1.449 στις εκατοντάδες είναι 1.400, γιατί το ψηφίο των δεκάδων είναι 4. Αν το κάνεις σε δύο βήματα, πρώτα στις δεκάδες (1.450) και μετά στις εκατοντάδες, βγάζεις 1.500, που είναι λάθος.
\item Όταν το ψηφίο που αυξάνεται είναι 9, η αύξηση «κουβαλάει»: $9.995\approx 10.000$ στις δεκάδες.
\item Όταν μετράς \textbf{από ως}, και τα δύο άκρα μετράνε. Από τη σελίδα 20 ως τη σελίδα 47 είναι $47-20+1=28$ σελίδες, όχι 27.
\end{itemize}
\end{prosoxi}

\tmima{Λυμένα παραδείγματα}

\paradeigma{Το ντουλάπι του Ηλία}
\tixero Το ντουλάπι του Ηλία ανοίγει με τριψήφιο κωδικό. Θυμάται ότι έχει τα ψηφία 3, 0 και 5, καθένα μία φορά, αλλά όχι με ποια σειρά.\par
\tithelo Όλους τους πιθανούς κωδικούς, σε αύξουσα σειρά.\par
\possindeo Δοκιμάζω όλες τις σειρές των τριών ψηφίων: 305, 350, 035, 053, 503, 530. Ο 035 και ο 053 δεν είναι τριψήφιοι αριθμοί, γιατί το 0 μπροστά δεν μετράει. Τους σβήνω.\par
\apantisi 305, 350, 503, 530.\par
\elegxos Είναι τέσσερις, και κανένας δεν ξεκινάει από 0. Ο μικρότερος είναι ο 305: ξεκινάει από το μικρότερο μη μηδενικό ψηφίο, το 3, και μετά βάζει το μικρότερο που μένει, το 0.

\paradeigma{Οι προβολές του Βασίλη}
\tixero Το βίντεο της σχολικής εκλογικής εκστρατείας του Βασίλη έχει 3.486.917 προβολές. (Ναι, ο Βασίλης υπερβάλλει, αλλά ας το δεχτούμε.)\par
\tithelo Να το στρογγυλοποιήσω (α) στις εκατοντάδες, (β) στις χιλιάδες, (γ) στις εκατοντάδες χιλιάδων.\par
\possindeo
(α) Τάξη: εκατοντάδες. Η αμέσως μικρότερη τάξη είναι οι δεκάδες, με ψηφίο 1. Επειδή $1<5$, το 9 μένει ίδιο και τα δεξιά γίνονται 0: $3.486.917\approx 3.486.900$.\par
(β) Τάξη: χιλιάδες, ψηφίο 6. Η εκατοντάδα έχει ψηφίο 9. Επειδή $9\ge 5$, το 6 γίνεται 7: $3.486.917\approx 3.487.000$.\par
(γ) Τάξη: εκατοντάδες χιλιάδων, ψηφίο 4. Οι δεκάδες χιλιάδων έχουν ψηφίο 8. Επειδή $8\ge 5$, το 4 γίνεται 5: $3.486.917\approx 3.500.000$.\par
\apantisi 3.486.900, 3.487.000 και 3.500.000.\par
\elegxos Στο (β) το 3.486.917 απέχει 83 από το 3.487.000 και 917 από το 3.486.000. Το πρώτο είναι πιο κοντά. Στο (γ) απέχει 13.083 από το 3.500.000 και 86.917 από το 3.400.000. Και πάλι το 3.500.000 είναι πιο κοντά.

\paradeigma{Πόσες σελίδες;}
\tixero Η Φανή διαβάζει από τη σελίδα 58 ως τη σελίδα 131 του βιβλίου της, και τις δύο μαζί.\par
\tithelo Πόσες σελίδες διαβάζει.\par
\possindeo Αν έβγαζα $131-58=73$, θα είχα αφήσει έξω την πρώτη σελίδα. Για να μετρήσω και τα δύο άκρα, προσθέτω 1: $131-58+1=74$.\par
\apantisi 74 σελίδες.\par
\elegxos Δοκιμή με μικρούς αριθμούς: από τη σελίδα 3 ως τη 5 είναι τρεις σελίδες (3, 4, 5). Ο τύπος δίνει $5-3+1=3$. Σωστά.

\tmima{Ασκήσεις}

\zoni{Το κατάλαβα}{εφαρμογή του κανόνα}

\ex Γράψε με ψηφία: α) τριακόσια δεκαοκτώ \quad β) εννιακόσια ένα \quad γ) δεκαπέντε χιλιάδες διακόσια τέσσερα \quad δ) δύο εκατομμύρια τριακόσιες χιλιάδες εβδομήντα.

\ex Γράψε με λέξεις: α) 64.307 \quad β) 1.050.009 \quad γ) 200.020.

\ex Ποιοι είναι οι τρεις προηγούμενοι αριθμοί του 1.000 και ποιοι οι δύο επόμενοι του 999.999; Ποιος είναι ο προηγούμενος του 0;

\ex Τα σκορ από το παιχνίδι του Βασίλη είναι 7.250, 7.205, 72.050, 7.520 και 5.270. Γράψε τα σε αύξουσα σειρά και μετά σε φθίνουσα.

\ex Βάλε το σωστό σύμβολο ($<$, $=$ ή $>$): α) $603\ \square\ 630$ \quad β) $9.099\ \square\ 9.100$ \quad γ) $45.012\ \square\ 45.021$ \quad δ) $100.000\ \square\ 99.999$ \quad ε) $7.008\ \square\ 7.008$ \quad στ) $12.345\ \square\ 1.234$.

\ex Σχεδίασε μια αριθμητική ευθεία με αρχή το Ο και μονάδα $\mathrm{ΟΑ}=3$ εκ. Τοποθέτησε τα σημεία Β, Γ, Δ και Ε σε αποστάσεις 9, 15, 21 και 27 εκ. από το Ο. Ποιοι αριθμοί αντιστοιχούν σε αυτά;

\ex Στρογγυλοποίησε στην πλησιέστερη εκατοντάδα: 248, 350, 1.249, 5.951, 19.999 και 70.049.

\ex Στρογγυλοποίησε τον αριθμό 2.846.593 (α) στις δεκάδες, (β) στις εκατοντάδες, (γ) στις χιλιάδες, (δ) στις δεκάδες χιλιάδων, (ε) στις εκατοντάδες χιλιάδων, (στ) στα εκατομμύρια.

\ex Σωστό ή λάθος;\par
α) Στον αριθμό 4.050.300 τα μηδενικά δείχνουν ότι δεν υπάρχουν εκατοντάδες χιλιάδων, χιλιάδες, δεκάδες και μονάδες.\quad β) Εκατό δεκάδες είναι χίλιες μονάδες.\quad γ) Από τον 40 ως τον 99 υπάρχουν 60 αριθμοί.\quad δ) Από τη σελίδα 20 ως τη σελίδα 29 είναι 9 σελίδες.\quad ε) Αν σήμερα είναι Δευτέρα, σε 15 ημέρες θα είναι πάλι Δευτέρα.\quad στ) Ανάμεσα στο 7 και στο 8 δεν υπάρχει άλλος φυσικός αριθμός.\quad ζ) Το 2.449 στρογγυλοποιημένο στις εκατοντάδες είναι 2.500.

\zoni{Το χρησιμοποιώ}{προβλήματα από τη ζωή}

\ex \textbf{Το λουκέτο της Φανής.} Το λουκέτο της ποδηλατοθήκης έχει τριψήφιο κωδικό με τα ψηφία 2, 0 και 8, καθένα μία φορά. Γράψε όλους τους δυνατούς κωδικούς από τον μικρότερο στον μεγαλύτερο. Πόσους θα δοκιμάσει η Φανή στη χειρότερη περίπτωση;

\ex \textbf{Το χωριό του παππού.} Η πινακίδα στην είσοδο του χωριού γράφει «Κάτοικοι: 1.846». Ο Νίκος λέει στους φίλους του «περίπου 2.000» και η Μυρτώ «περίπου 1.800». Σε ποια τάξη στρογγυλοποίησε ο καθένας; Ποιος είναι πιο κοντά; Αν ο δήμαρχος πρέπει να τυπώσει ψηφοδέλτια για όλους τους κατοίκους, μπορεί να χρησιμοποιήσει το «περίπου»; Γιατί;

\ex \textbf{Τα βήματα του Νίκου.} Η εφαρμογή στο κινητό του Νίκου γράφει ότι σήμερα έκανε 14.992 βήματα. Ο Νίκος λέει «Στρογγυλά, 15.000, πέτυχα τον στόχο!». Ο στόχος είναι τουλάχιστον 15.000 βήματα. Έφτασε τον στόχο; Πόσα του λείπουν; Γιατί εδώ η στρογγυλοποίηση δεν βοηθάει;

\ex \textbf{Τα πλακάκια του θείου.} Ο θείος Λευτέρης χρειάζεται 2.314 πλακάκια για τον κήπο. Πωλούνται σε κιβώτια των 100. Ο Στέλιος στρογγυλοποιεί στην πλησιέστερη εκατοντάδα και λέει «23 κιβώτια». Φτάνουν; Πόσα κιβώτια πρέπει να αγοράσει και πόσα πλακάκια θα περισσέψουν;

\ex \textbf{Το βιβλίο της Μυρτώς.} Η Μυρτώ διαβάζει από τη σελίδα 204 ως τη σελίδα 311, και τις δύο μαζί. Πόσες σελίδες διαβάζει; Αν διαβάζει 9 σελίδες την ημέρα, σε πόσες ημέρες θα τελειώσει;

\ex \textbf{Ο φράχτης.} Ο παππούς φράζει ένα κομμάτι του κήπου μήκους 36 μ. Βάζει έναν πάσσαλο κάθε 3 μ., και έναν στην αρχή και έναν στο τέλος. Πόσους πασσάλους θα χρειαστεί;

\ex \textbf{Οι κάρτες του αγώνα.} Οι κάρτες του σχολικού αγώνα είναι αριθμημένες από το 501 ως το 780. Πόσες κάρτες υπάρχουν; Πόσες από αυτές έχουν περιττό αριθμό;

\ex \textbf{Ο άξονας του Βασίλη.} Ο Βασίλης σχεδιάζει έναν άξονα για τις βαθμολογίες του παιχνιδιού του. Τα σημεία 0 και 1.000 απέχουν 4 εκ. (δες το σχήμα). Ποιοι αριθμοί αντιστοιχούν στα σημεία Α, Β, Γ και Δ;

\begin{center}
\begin{tikzpicture}[x=1.15cm]
\draw[-{Stealth}] (-0.3,0)--(9.6,0);
\draw (0,-0.12)--(0,0.12) node[above,font=\small] {0};
\draw (4,-0.12)--(4,0.12) node[above,font=\small] {1.000};
\foreach \x/\n in {1/Α,3/Β,5/Γ,7.5/Δ}{\draw[gblue,thick] (\x,-0.15)--(\x,0.15); \node[below,font=\small\bfseries,gblue] at (\x,-0.15) {\n};}
\end{tikzpicture}
\end{center}

\zoni{Το σκέφτομαι}{μια ερώτηση που πάει πιο μακριά}

\ex Πόσες φορές γράφεται το ψηφίο 3, αν γράψεις όλους τους αριθμούς από το 1 ως το 200;

\ex Ποιος είναι ο μικρότερος και ποιος ο μεγαλύτερος φυσικός αριθμός που, στρογγυλοποιημένος στις χιλιάδες, δίνει 8.000;

\ex Ποιος είναι ο μεγαλύτερος και ποιος ο μικρότερος πενταψήφιος αριθμός που έχει όλα τα ψηφία του διαφορετικά;

\ex Ο Νίκος λέει: «Θα βρω τον μεγαλύτερο φυσικό αριθμό και θα τον γράψω στον πίνακα». Η Ζωή λέει: «Δεν υπάρχει». Εξήγησε με λόγια γιατί έχει δίκιο η Ζωή.

\ex Ο Ηλίας δοκιμάζει στο κινητό: $7+5=12$, $3+9=12$. Ρωτάει την Πολυξένη: «Αν προσθέσω δύο φυσικούς, βγαίνει πάντα φυσικός. Αν αφαιρέσω;» Η Πολυξένη γράφει $3-5$. Μπορεί να υπάρχει αυτό το αποτέλεσμα στους φυσικούς; Τι νομίζεις ότι θα χρειαστούμε;

\zoni{Για το σπίτι}{πιο μακριά, με χαρτί και χάρακα}

\ex \textbf{Το θερμόμετρο της Αγγελικής.} Η Αγγελική φτιάχνει ένα χάρτινο θερμόμετρο για την τάξη. Σημειώνει το $0^\circ$ και το $40^\circ$ και η απόσταση ανάμεσά τους είναι 20 εκ. Πού θα βάλει τα σημάδια για τους $10^\circ$, $20^\circ$, $25^\circ$ και $35^\circ$;

\ex Δοκίμασε τις αφαιρέσεις $9-3$, $15-7$ και $21-5$. Τι είδους αριθμούς βρίσκεις; Εξήγησε γιατί η διαφορά δύο περιττών αριθμών είναι πάντα άρτιος.

\ex Πόσοι τριψήφιοι αριθμοί υπάρχουν; (Βοήθεια: ο μικρότερος είναι ο 100 και ο μεγαλύτερος ο 999, και τα δύο άκρα μετράνε.)

\paragrafos{Ευκλείδεια διαίρεση: το υπόλοιπο που μένει}{Α.1.2, «Ταυτότητα Ευκλείδειας Διαίρεσης»}

\begin{anagke}
Το σχολείο οργανώνει την πρώτη του εκδρομή. Θα πάνε και τα 387 παιδιά, με λεωφορεία των 49 θέσεων. Η κυρία Αλεξάνδρα ζητάει από την τάξη να βρει πόσα λεωφορεία χρειάζονται.

Ο Στέλιος, που αγαπάει το κομπιουτεράκι, κάνει $387:49$ και η οθόνη γράφει $7{,}897959\dots$ «Εφτά κόμμα εννιά!» λέει περήφανος. «Άρα εφτά λεωφορεία». Η Ρήνη τον κοιτάει. «Και τα υπόλοιπα παιδιά; Μένουν στην αυλή;» Ο Στέλιος ξαναδιαβάζει την οθόνη. «Μα το 0,897 τι είναι; Δεν υπάρχει 0,897 λεωφορείου». Η Ζωή κάνει τη διαίρεση στο χαρτί, όπως τη μάθαμε: πηλίκο 7 και κάτι περισσεύει. «Αυτό που περισσεύει είναι παιδιά, όχι κόμματα. Εκεί βρίσκεται η απάντηση».

Σε πολλά προβλήματα δεν ενδιαφέρει μόνο πόσες φορές χωράει ο ένας αριθμός στον άλλον, ενδιαφέρει και τι περισσεύει. Σε αυτή την παράγραφο μαθαίνουμε να γράφουμε τη διαίρεση με μία ισότητα, την ταυτότητα της Ευκλείδειας διαίρεσης, και να τη χρησιμοποιούμε για να λύνουμε προβλήματα.
\end{anagke}

\begin{lexi}{πηλίκο, υπόλοιπο, Ευκλείδεια}{πηλίκο, υπόλοιπο}
Το \lexiword{πηλίκο} έρχεται από το αρχαίο «πηλίκος», που σημαίνει «πόσο μεγάλος»: η διαίρεση απαντά στο «πόσο μεγάλο μερίδιο παίρνει ο καθένας». Το \lexiword{υπόλοιπο} είναι κυριολεκτικά αυτό που μένει, το «λοιπό» που δεν μοιράστηκε.

Η \lexiword{Ευκλείδεια διαίρεση} πήρε το όνομά της από τον Ευκλείδη, τον Έλληνα μαθηματικό που έζησε στην Αλεξάνδρεια γύρω στο 300 π.Χ. και έγραψε τα «Στοιχεία», το βιβλίο όπου συγκέντρωσε τα μαθηματικά της εποχής του.
\end{lexi}

\begin{metafrasi}
Μοιράζω 43 μπάλες σε ομάδες των 7: γίνονται 6 ομάδες και μένει 1 μπάλα. & $43=7\cdot 6+1$\\
Δεν περισσεύει τίποτα (τέλεια διαίρεση). & $\Delta=\delta\cdot\pi$, υπόλοιπο $0$\\
Το υπόλοιπο είναι πάντα μικρότερο από τον διαιρέτη. & $\upsilon<\delta$\\
Στη διαίρεση με το 8 τα πιθανά υπόλοιπα. & $0,1,2,3,4,5,6,7$\\
Δεν διαιρούμε ποτέ με το μηδέν. & $\delta\neq 0$\\
\end{metafrasi}

\tmima{Η θεωρία}

\begin{kanonas}{ταυτότητα της Ευκλείδειας διαίρεσης}
Αν $\Delta$ και $\delta$ είναι δύο φυσικοί αριθμοί με $\delta\neq 0$, τότε υπάρχουν δύο άλλοι φυσικοί $\pi$ και $\upsilon$ έτσι ώστε
$$\Delta=\delta\cdot\pi+\upsilon,\qquad\text{με}\quad \upsilon<\delta.$$
\end{kanonas}

Ο $\Delta$ λέγεται \textbf{διαιρετέος}, ο $\delta$ \textbf{διαιρέτης}, ο $\pi$ \textbf{πηλίκο} και ο $\upsilon$ \textbf{υπόλοιπο}. Για παράδειγμα, από τη διαίρεση $43:7$ έχουμε $\pi=6$ και $\upsilon=1$, δηλαδή $43=7\cdot 6+1$ με $1<7$.

Αν το υπόλοιπο είναι $0$, η διαίρεση λέγεται \textbf{τέλεια} και γράφεται $\Delta=\delta\cdot\pi$. Τότε ισχύει και $\Delta:\delta=\pi$ και $\Delta:\pi=\delta$. Η τέλεια διαίρεση είναι η πράξη αντίστροφη του πολλαπλασιασμού.

\begin{idiotita}{ειδικές περιπτώσεις}
\begin{itemize}[leftmargin=1.2em,itemsep=0.3ex,topsep=2pt]
\item Ο διαιρέτης δεν μπορεί να είναι $0$: $\delta\neq 0$.
\item $\alpha:\alpha=1$ (όταν $\Delta=\delta$ το πηλίκο είναι $1$).
\item $\alpha:1=\alpha$ (όταν $\delta=1$ το πηλίκο είναι ο διαιρετέος).
\item $0:\alpha=0$ (όταν $\Delta=0$ το πηλίκο είναι $0$).
\end{itemize}
\end{idiotita}

\begin{prosoxi}
Μια ισότητα της μορφής $\Delta=\delta\cdot\pi+\upsilon$ δεν είναι πάντα Ευκλείδεια διαίρεση. Πρέπει να ισχύει και $\upsilon<\delta$. Για παράδειγμα η $50=7\cdot 6+8$ είναι αληθινή ισότητα, αλλά δεν είναι Ευκλείδεια διαίρεση, γιατί $8>7$: στο $50$ το $7$ χωράει ακόμη μία φορά ($50=7\cdot 7+1$).
\end{prosoxi}

\tmimasmall{Τι υπόλοιπα μπορούν να βγουν}

Επειδή $\upsilon<\delta$, στη διαίρεση με το $5$ τα πιθανά υπόλοιπα είναι μόνο $0,1,2,3,4$. Στη διαίρεση με το $\delta$ υπάρχουν ακριβώς $\delta$ διαφορετικά υπόλοιπα, από το $0$ ως το $\delta-1$. Γι' αυτό τα προβλήματα που επαναλαμβάνονται κυκλικά (ημέρες της εβδομάδας, ώρες του ρολογιού) λύνονται με Ευκλείδεια διαίρεση.

\tmima{Λυμένα παραδείγματα}

\paradeigma{Ποιες ισότητες είναι Ευκλείδεια διαίρεση;}
\tixero Οι ισότητες α) $54=25\cdot 3+9$, β) $370=46\cdot 8+2$, γ) $680=51\cdot 12+68$.\par
\tithelo Σε ποιες ισχύει $\upsilon<\delta$ και ποιοι είναι οι $\Delta,\delta,\pi,\upsilon$;\par
\possindeo Στην α) είναι $9<25$: Ευκλείδεια διαίρεση με $\Delta=54$, $\delta=25$, $\pi=3$, $\upsilon=9$. Στη β) το $2$ είναι μικρότερο και από τον $46$ και από τον $8$, άρα μπορούμε να πούμε $\delta=46,\pi=8$ ή $\delta=8,\pi=46$ (και στις δύο $\upsilon=2$). Στη γ) το $68$ είναι μεγαλύτερο και από τον $51$ και από τον $12$, άρα δεν είναι διαίρεση.\par
\apantisi Οι α) και β) εκφράζουν Ευκλείδεια διαίρεση, η γ) όχι.

\paradeigma{Οι καρέκλες της αίθουσας}
\tixero Κάθε αίθουσα πρέπει να έχει $12$ θρανία. Σήμερα ήρθαν $3$ φορτηγά που το καθένα μετέφερε $35$ θρανία.\par
\tithelo Πόσες αίθουσες συμπληρώνονται και πόσα θρανία χρειάζονται ακόμη για να συμπληρωθεί μία ακόμη;\par
\possindeo Συνολικά $3\cdot 35=105$ θρανία. Κάνω τη διαίρεση $105:12$: $105=12\cdot 8+9$. Άρα $8$ αίθουσες συμπληρώνονται και περισσεύουν $9$ θρανία. Για την επόμενη χρειάζονται ακόμη $12-9=3$.\par
\apantisi $8$ αίθουσες, και χρειάζονται άλλα $3$ θρανία.\par
\elegxos $12\cdot 8+9=96+9=105$ και $9<12$.

\paradeigma{Η μέρα της εβδομάδας}
\tixero Σήμερα είναι Δευτέρα.\par
\tithelo Τι μέρα θα είναι σε $100$ ημέρες;\par
\possindeo Οι ημέρες επαναλαμβάνονται ανά $7$. Κάνω τη διαίρεση: $100=7\cdot 14+2$. Μετά από $14$ πλήρεις εβδομάδες ξαναγυρνάμε σε Δευτέρα, και μένουν $2$ ημέρες: Τρίτη, Τετάρτη.\par
\apantisi Τετάρτη.\par
\elegxos $7\cdot 14=98$ ημέρες είναι Δευτέρα, και $2$ ακόμη φέρνουν Τετάρτη.

\tmima{Ασκήσεις}

\zoni{Το κατάλαβα}{εφαρμογή της ταυτότητας}

\ex Συμπλήρωσε: α) Στη διαίρεση $47:6$ ο διαιρετέος είναι \dots, ο διαιρέτης \dots, το πηλίκο \dots και το υπόλοιπο \dots β) Ο διαιρέτης δεν μπορεί να είναι \dots γ) Αν ο διαιρετέος ισούται με τον διαιρέτη, το πηλίκο είναι \dots δ) Αν ο διαιρετέος είναι $0$, το πηλίκο είναι \dots

\ex Ποιες από τις ισότητες εκφράζουν Ευκλείδεια διαίρεση; α) $47=6\cdot 7+5$ \quad β) $50=7\cdot 6+8$ \quad γ) $29=4\cdot 7+1$ \quad δ) $100=9\cdot 10+10$ \quad ε) $63=9\cdot 7+0$.

\ex Κάνε τις διαιρέσεις και τις δοκιμές τους: α) $135:5$ \quad β) $95:4$ \quad γ) $276:12$ \quad δ) $1.000:7$. Ποιες είναι τέλειες;

\ex Ποια υπόλοιπα μπορεί να έχει μια διαίρεση με διαιρέτη το $5$; Ποια με διαιρέτη το $8$; Ποια με διαιρέτη το $1$;

\ex Σωστό ή λάθος; α) Η διαίρεση $452:2$ είναι τέλεια. \quad β) Ο διαιρέτης μπορεί να είναι $0$. \quad γ) Το υπόλοιπο είναι πάντα μικρότερο από τον διαιρέτη. \quad δ) Αν ο διαιρέτης είναι $1$, το πηλίκο ισούται με τον διαιρετέο. \quad ε) $0:7=7$. \quad στ) $\alpha:\alpha=\alpha$.

\ex Γράψε την ταυτότητα της Ευκλείδειας διαίρεσης για: α) $83:9$ \quad β) $250:25$ \quad γ) $7:10$ \quad δ) $0:5$.

\zoni{Το χρησιμοποιώ}{προβλήματα από τη ζωή}

\ex \textbf{Τα βαρέλια.} Ο παππούς Γιάννης έχει $350$ λίτρα κρασί και βάζει $12$ λίτρα σε κάθε βαρέλι. Πόσα βαρέλια γεμίζουν και πόσα λίτρα περισσεύουν;

\ex \textbf{Η εκδρομή.} Στην εκδρομή θέλουν να πάνε $250$ παιδιά και κάθε λεωφορείο έχει $45$ θέσεις. Πόσα λεωφορεία χρειάζονται; Πόσες θέσεις μένουν άδειες στο τελευταίο;

\ex \textbf{Οι ημέρες.} Σήμερα είναι Τετάρτη. Τι μέρα θα είναι μετά από $45$ ημέρες; Και μετά από $100$ ημέρες;

\ex \textbf{Τα βραχιόλια.} Η Ρήνη έχει $100$ χάντρες και φτιάχνει βραχιόλια με $8$ χάντρες το καθένα. Πόσα βραχιόλια φτιάχνει και πόσες χάντρες περισσεύουν; Πόσες χάντρες της λείπουν για να φτιάξει ένα ακόμη;

\ex \textbf{Το όνομα που επαναλαμβάνεται.} Η Αγγελική γράφει συνέχεια τη λέξη $\mathrm{ΑΓΓΕΛΙΚΗ}$ σε μια λωρίδα χαρτί: $\mathrm{ΑΓΓΕΛΙΚΗΑΓΓΕΛΙΚΗ}\dots$ Ποιο γράμμα βρίσκεται στην $100$ή θέση;

\zoni{Το σκέφτομαι}{ερωτήσεις με λίγη περισσότερη σκέψη}

\ex Ένας αριθμός διαιρούμενος με το $7$ δίνει πηλίκο $12$ και υπόλοιπο $6$. Ποιος είναι ο αριθμός;

\ex Ένας αριθμός διαιρείται τέλεια με το $14$ και δίνει πηλίκο $201$. Ποιος είναι ο αριθμός;

\ex Ένας αριθμός διαιρούμενος με το $9$ έχει πηλίκο $5$ και το μεγαλύτερο δυνατό υπόλοιπο. Ποιος είναι; Τι πηλίκο και τι υπόλοιπο δίνει ο επόμενός του;

\ex Ο Βασίλης λέει: «Αν ένας αριθμός διαιρούμενος με το $4$ δίνει υπόλοιπο $3$, τότε ο διπλάσιός του διαιρούμενος με το $4$ θα δώσει υπόλοιπο $6$». Εξήγησε γιατί αυτό δεν μπορεί να ισχύει και βρες το σωστό υπόλοιπο. (Δοκίμασε με τον αριθμό $7$.)

\zoni{Για το σπίτι}{εξάσκηση}

\ex Υπολόγισε τις τιμές: α) $35-12:3+3\cdot 3$ \quad β) $63:9+63:7-12:2$ \quad γ) $(5+12:3):3-24:(2\cdot 2+16:4)$.

\ex \textbf{Οι συνοδοί.} Στο καταφύγιο ζώων της περιοχής πρέπει να υπάρχει ένας εθελοντής ανά $15$ σκυλάκια, το πολύ. Αν στο καταφύγιο υπάρχουν $162$ σκυλάκια, ποιος είναι ο μικρότερος αριθμός εθελοντών που χρειάζεται;

\paragrafos{Δυνάμεις φυσικών αριθμών: ποιος πάει πρώτος}{Α.1.3, «Δυνάμεις φυσικών αριθμών»}

\begin{anagke}
Παρασκευή απόγευμα. Η Λίνα φτιάχνει κωδικό για το καινούργιο της κινητό: τέσσερα ψηφία. «Πόσοι διαφορετικοί κωδικοί υπάρχουν;» ρωτάει η Ζωή. Η Λίνα σκέφτεται δυνατά: «Για το πρώτο ψηφίο έχω 10 επιλογές, από το 0 ως το 9. Για το δεύτερο πάλι 10, και για το τρίτο, και για το τέταρτο». Ο Ηλίας γράφει $10\cdot 10\cdot 10\cdot 10$ και βρίσκει 10.000. «Και αν ο κωδικός είχε είκοσι ψηφία;» ρωτάει. Γράφει είκοσι δεκάρια στη σειρά και χάνει το μέτρημα στη μέση.

Ο Άρης, που βιάζεται πάντα, γράφει σε ένα χαρτάκι $2+3\cdot 4$ και ρωτάει τη Ζωή πόσο κάνει. «Δεκατέσσερα», λέει εκείνη. «Είκοσι!» φωνάζει ο Άρης. «Δύο και τρία πέντε, επί τέσσερα, είκοσι». Ο Θοδωρής, που περνάει από δίπλα, τους κοιτάει και λέει μόνο: «Και οι δύο έχετε ένα κομπιουτεράκι στην τσέπη. Δείτε τι βγάζει». Το κομπιουτεράκι δίνει 14.

Δύο προβλήματα. Το πρώτο: πώς γράφεται σύντομα ο ίδιος παράγοντας πολλές φορές; Το δεύτερο: όταν μια παράσταση έχει πολλές πράξεις, ποια γίνεται πρώτη, και ποιος το αποφασίζει;
\end{anagke}

\begin{lexi}{δύναμη, εκθέτης, βάση}{δύναμη, εκθέτης}
Η λέξη \lexiword{δύναμη} στα μαθηματικά δεν έχει σχέση με τα μπράτσα. Ο αριθμός «έχει τη δύναμη» να πολλαπλασιάζεται με τον εαυτό του πολλές φορές. Το $2^{10}$ διαβάζεται «το δύο στη δέκατη».

Ο \lexiword{εκθέτης} είναι το μικρό νούμερο πάνω δεξιά. Το «εκθέτω» σημαίνει «δείχνω»: ο εκθέτης δείχνει πόσες φορές μπαίνει ο παράγοντας. Ο αριθμός που πολλαπλασιάζεται λέγεται \lexiword{βάση}, γιατί πάνω του «στηρίζεται» το σύμβολο.
\end{lexi}

\begin{metafrasi}
Το 2 πολλαπλασιασμένο με τον εαυτό του 10 φορές. & $2^{10}=2\cdot 2\cdot\dots\cdot 2$ (10 παράγοντες)\\
Το τετράγωνο του 5. & $5^2=5\cdot 5=25$\\
Ο κύβος του 3. & $3^3=3\cdot 3\cdot 3=27$\\
Το 7 στην πρώτη. & $7^1=7$\\
Ένα εκατομμύριο, με δύναμη του 10. & $10^6=1.000.000$\\
Πρώτα οι δυνάμεις, μετά οι πολλαπλασιασμοί. & $2+3\cdot 4^2=2+3\cdot 16=50$\\
\end{metafrasi}

\tmima{Η θεωρία}

\begin{orismos}{δύναμη}
Το γινόμενο $\alpha\cdot\alpha\cdot\alpha\cdots\alpha$ με $\nu$ παράγοντες ίσους με το $\alpha$ λέγεται \textbf{δύναμη} του $\alpha$ στη $\nu$ και γράφεται $\alpha^{\nu}$. Ο $\alpha$ είναι η \textbf{βάση} και ο $\nu$ ο \textbf{εκθέτης}.
\end{orismos}

Η δεύτερη δύναμη $\alpha^2$ λέγεται και \textbf{τετράγωνο} του $\alpha$ (είναι το εμβαδό τετραγώνου με πλευρά $\alpha$). Η τρίτη δύναμη $\alpha^3$ λέγεται και \textbf{κύβος} του $\alpha$ (είναι ο όγκος κύβου με ακμή $\alpha$). Η πρώτη δύναμη $\alpha^1$ είναι ο ίδιος ο αριθμός $\alpha$. Οι δυνάμεις του 1 είναι όλες 1: $1^{\nu}=1$.

\begin{prosoxi}
Το $2^3$ δεν είναι $2\cdot 3=6$ αλλά $2\cdot 2\cdot 2=8$. Ο εκθέτης δείχνει πόσες φορές γράφεται η βάση, όχι επί πόσο την πολλαπλασιάζουμε.
\end{prosoxi}

\tmimasmall{Οι δυνάμεις του 10}

Οι δυνάμεις του 10 έχουν τόσα μηδενικά όσος ο εκθέτης:
$$10^2=100,\qquad 10^3=1.000,\qquad 10^6=1.000.000.$$
Γι' αυτό η δεκαδική μας γραφή αναλύεται εύκολα. Για παράδειγμα:
$$7.604=7\cdot 10^3+6\cdot 10^2+0\cdot 10+4.$$

\tmimasmall{Αριθμητική παράσταση και προτεραιότητα}

Μια \textbf{αριθμητική παράσταση} είναι μια σειρά αριθμών που συνδέονται με σύμβολα πράξεων. Ο αριθμός που βγαίνει όταν κάνουμε όλες τις πράξεις λέγεται \textbf{τιμή} της. Για να βγάζουμε όλοι την ίδια τιμή, όλος ο κόσμος έχει συμφωνήσει σε μια σειρά.

\begin{kanonas}{προτεραιότητα των πράξεων}
\begin{enumerate}[leftmargin=1.6em,itemsep=0.2ex,topsep=2pt]
\item Πρώτα οι πράξεις μέσα στις \textbf{παρενθέσεις}, με την ίδια σειρά που θα βρεις παρακάτω.
\item Μετά οι \textbf{δυνάμεις}.
\item Μετά οι \textbf{πολλαπλασιασμοί και οι διαιρέσεις}, από αριστερά προς τα δεξιά.
\item Τελευταίες οι \textbf{προσθέσεις και οι αφαιρέσεις}, από αριστερά προς τα δεξιά.
\end{enumerate}
\end{kanonas}

Η συμφωνία αυτή δεν είναι «φυσικός νόμος». Είναι συμφωνία, όπως ο κανόνας ότι στον δρόμο πάμε δεξιά. Αν δεν υπήρχε, το $2+3\cdot 4$ θα είχε δύο απαντήσεις (14 και 20) και δύο μηχανικοί δεν θα συμφωνούσαν ποτέ σε έναν υπολογισμό. Με τη συμφωνία η απάντηση είναι μία, 14, και η Ζωή έχει δίκιο. Αν θέλουμε το 20, γράφουμε $(2+3)\cdot 4$: η παρένθεση λέει ρητά «αυτό πρώτα».

\begin{prosoxi}
\begin{itemize}[leftmargin=1.2em,itemsep=0.3ex,topsep=2pt]
\item Στο $2\cdot 3^2$ ο εκθέτης αφορά μόνο το 3: $2\cdot 9=18$. Για να τετραγωνίσεις όλο το γινόμενο, γράφεις $(2\cdot 3)^2=36$.
\item Οι πράξεις του ίδιου επιπέδου γίνονται από αριστερά προς τα δεξιά: $24:4:2=6:2=3$, όχι $24:(4:2)=12$.
\end{itemize}
\end{prosoxi}


\tmimasmall{Ιστορική αναδρομή: πώς έγραφαν τους αριθμούς}

Το δεκαδικό σύστημα, που σου φαίνεται αυτονόητο, χρειάστηκε χιλιάδες χρόνια για να διαμορφωθεί. Οι Σουμέριοι και οι Βαβυλώνιοι (3η–2η χιλιετία π.Χ.) έγραφαν τους αριθμούς σε πήλινες πινακίδες. Οι Αιγύπτιοι είχαν ένα ιερογλυφικό για κάθε τάξη: 1, 10, 100, 1.000, \dots, 1.000.000. Οι Ίωνες, τον 5ο αιώνα π.Χ., έφτιαξαν το \textbf{αλφαβητικό σύστημα}: κάθε γράμμα του αλφαβήτου σήμαινε έναν αριθμό.

\begin{center}
\small\renewcommand{\arraystretch}{1.25}\setlength{\tabcolsep}{4pt}
\begin{tabular}{cccccccccc}
α & β & γ & δ & ε & \arch{ϛ} & ζ & η & θ & \\
1 & 2 & 3 & 4 & 5 & 6 & 7 & 8 & 9 & \\
ι & κ & λ & μ & ν & ξ & ο & π & \arch{ϙ} & \\
10 & 20 & 30 & 40 & 50 & 60 & 70 & 80 & 90 & \\
ρ & σ & τ & υ & φ & χ & ψ & ω & \arch{ϡ} & \\
100 & 200 & 300 & 400 & 500 & 600 & 700 & 800 & 900 & \\
\end{tabular}
\end{center}

Έτσι το 831 γράφεται ωλα ($800+30+1$). Οι \textbf{Ρωμαίοι} είχαν σύμβολα για τους αριθμούς $1,\,5,\,10,\,50,\,100,\,500,\,1.000$:

\begin{center}
\small\setlength{\tabcolsep}{6pt}
\begin{tabular}{ccccccc}
I & V & X & L & C & D & M\\
1 & 5 & 10 & 50 & 100 & 500 & 1.000\\
\end{tabular}
\end{center}

Πρόσθεταν από αριστερά προς τα δεξιά (II = 2, XVI = 16) και αφαιρούσαν όταν ένα μικρό σύμβολο προηγείται: IV = 4 ($5-1$), IX = 9, XL = 40, CM = 900. Τα σύμβολα που χρησιμοποιούμε σήμερα (0 ως 9) προέρχονται από τους Ινδούς, έφτασαν στην Ευρώπη μέσω των Αράβων και γι' αυτό λέγονται \textbf{ινδοαραβικά}. Εμφανίζονται πρώτη φορά στο έργο του Αλ-Χουαρίζμι (9ος αιώνας μ.Χ.).

Δεν είναι ανάγκη να έχουμε δέκα ψηφία. Ο υπολογιστής χρησιμοποιεί μόνο το 0 και το 1, στο \textbf{δυαδικό σύστημα}, όπου αντί για μονάδες, δεκάδες, εκατοντάδες έχουμε μονάδες, δυάδες, τετράδες, οκτάδες, δεκαεξάδες \dots

\begin{center}
\small\setlength{\tabcolsep}{4pt}
\begin{tabular}{l|cccccccccccc}
Δεκαδικό & 0 & 1 & 2 & 3 & 4 & 5 & 6 & 7 & 8 & 9 & 10 & 11\\
Δυαδικό & 0 & 1 & 10 & 11 & 100 & 101 & 110 & 111 & 1000 & 1001 & 1010 & 1011\\
\end{tabular}
\end{center}

Οι Πυθαγόρειοι πίστευαν ότι «τα πάντα είναι αριθμός». Έβαζαν τους αριθμούς σε σχήματα: οι \textbf{τρίγωνοι} αριθμοί $1,3,6,10,\dots$ φτιάχνουν τρίγωνα από τελείες και οι \textbf{τετράγωνοι} $1,4,9,16,\dots$ τετράγωνα. Οι τετράγωνοι είναι οι δεύτερες δυνάμεις.
\tmima{Λυμένα παραδείγματα}

\paradeigma{Οι κωδικοί της Λίνας}
\tixero Ο κωδικός έχει 4 ψηφία και κάθε ψηφίο έχει 10 επιλογές (0 ως 9).\par
\tithelo Πόσοι είναι οι δυνατοί κωδικοί, και πώς γράφεται σύντομα ο υπολογισμός;\par
\possindeo Για κάθε ψηφίο έχω 10 επιλογές, και οι επιλογές συνδυάζονται, άρα πολλαπλασιάζω: $10\cdot 10\cdot 10\cdot 10$. Είναι τέσσερα δεκάρια, δηλαδή $10^4$. Υπολογίζω βήμα βήμα: $10\cdot 10=100$, $100\cdot 10=1.000$, $1.000\cdot 10=10.000$.\par
\apantisi $10^4=10.000$ κωδικοί.\par
\elegxos Οι κωδικοί είναι οι αριθμοί από το 0000 ως το 9999. Από την παράγραφο 1.1: $9.999-0+1=10.000$. Βγήκε το ίδιο.

\paradeigma{Η παράσταση με πολλές πράξεις}
\tixero Η παράσταση $(2\cdot 5)^4+4\cdot(3+2)^2$.\par
\tithelo Την τιμή της.\par
\possindeo Πρώτα οι παρενθέσεις: $2\cdot 5=10$ και $3+2=5$. Έχω $10^4+4\cdot 5^2$. Μετά οι δυνάμεις: $10^4=10.000$ και $5^2=25$. Μετά ο πολλαπλασιασμός: $4\cdot 25=100$. Τέλος η πρόσθεση: $10.000+100=10.100$.\par
\apantisi $10.100$.\par
\elegxos Το $10^4$ έχει 4 μηδενικά, 10.000. Το άθροισμα με το 100 δίνει 10.100, όχι κάτι πιο μικρό από το 10.000. Σωστά.

\paradeigma{Ο αριθμός του Ηλία}
\tixero Ο Ηλίας γράφει ένα αριθμό σε ανάπτυγμα: $5\cdot 10^4+3\cdot 10^2+8$.\par
\tithelo Ποιος αριθμός είναι;\par
\possindeo $5\cdot 10^4=5\cdot 10.000=50.000$. $3\cdot 10^2=300$. Άρα $50.000+300+8=50.308$.\par
\apantisi $50.308$.\par
\elegxos Λείπουν οι χιλιάδες και οι δεκάδες, που σημαίνουν 0 στη θέση τους: 50.308. Τα μηδενικά έχουν μπει σωστά.

\tmima{Ασκήσεις}

\zoni{Το κατάλαβα}{εφαρμογή του κανόνα}

\ex Γράψε ως δύναμη: α) $5\cdot 5\cdot 5$ \quad β) $2\cdot 2\cdot 2\cdot 2\cdot 2$ \quad γ) $10\cdot 10\cdot 10\cdot 10$ \quad δ) $7\cdot 7$ \quad ε) $9$.

\ex Γράψε ως γινόμενο και υπολόγισε: α) $2^4$ \quad β) $3^3$ \quad γ) $5^2$ \quad δ) $10^5$ \quad ε) $1^{100}$ \quad στ) $6^1$.

\ex Ποια είναι η βάση και ποιος ο εκθέτης στα $4^5$, $12^2$, $1^{9}$ και $10^8$;

\ex Υπολόγισε το τετράγωνο και τον κύβο των αριθμών 2, 3, 4, 5 και 10. Φτιάξε πίνακα.

\ex Υπολόγισε την τιμή: α) $12+3\cdot 4$ \quad β) $(12+3)\cdot 4$ \quad γ) $24:4:2$ \quad δ) $100-5\cdot 4$ \quad ε) $2+3^2$ \quad στ) $(2+3)^2$.

\ex Βρες την τιμή: α) $(2\cdot 5)^4+4\cdot(3+2)^2$ \quad β) $(2+3)^3-8\cdot 3^2$ \quad γ) $3\cdot(4+2)^2-10\cdot 3$.

\ex Γράψε το ανάπτυγμα με δυνάμεις του 10: α) 7.604 \quad β) 53.280 \quad γ) 400.065.

\ex Γράψε με ψηφία τον αριθμό: α) $3\cdot 10^3+2\cdot 10^2+5$ \quad β) $9\cdot 10^5+4\cdot 10$ \quad γ) $10^6+10^2$.

\ex Σωστό ή λάθος; α) $2^3=6$ \quad β) $10^4=10.000$ \quad γ) $5^2=2\cdot 5$ \quad δ) $1^{57}=1$ \quad ε) $3+2\cdot 5=25$ \quad στ) $(1+1)^2=1^2+1^2$.

\zoni{Το χρησιμοποιώ}{προβλήματα από τη ζωή}

\ex \textbf{Το λουκέτο.} Ένα λουκέτο έχει τροχούς με ψηφία από το 0 ως το 9. Πόσους διαφορετικούς κωδικούς έχει αν έχει 3 τροχούς; Αν έχει 4; Αν έχει 5; Γράψε κάθε απάντηση και ως δύναμη του 10.

\ex \textbf{Οι διακόπτες.} Στην αποθήκη του θείου Λευτέρη ο πίνακας έχει διακόπτες που ο καθένας έχει δύο θέσεις: ανοιχτός ή κλειστός. Πόσους διαφορετικούς συνδυασμούς θέσεων έχουν 3 διακόπτες; Πόσους 5; Πόσους 10; Γράψε κάθε απάντηση και ως δύναμη του 2.

\ex \textbf{Ο κύβος του θείου.} Ο θείος Λευτέρης έχει δοχείο σε σχήμα κύβου με ακμή 4 δεκατόμετρα. Πόσα κυβικά δεκατόμετρα χωράει; (Θυμήσου ότι ο όγκος κύβου είναι $\alpha^3$.)

\ex \textbf{Ο κήπος.} Ο τετράγωνος κήπος της γιαγιάς Ρήνης έχει πλευρά 12 μ. Βρες το εμβαδό του. Αν διπλασιαστεί η πλευρά, ο κήπος γίνεται διπλάσιος σε εμβαδό; Υπολόγισε.

\ex \textbf{Τα σκακιστικά τετράγωνα.} Η σκακιέρα έχει 8 γραμμές και 8 στήλες. Πόσα τετράγωνα έχει; Γράψε τον υπολογισμό με δύναμη.

\ex \textbf{Ο λογαριασμός του Άρη.} Ο Άρης αγοράζει 3 μπάλες των 8 € και ένα φανελάκι των 12 €. Γράφει $3\cdot 8+12$ και βρίσκει 36. Η Μυρτώ γράφει $3\cdot(8+12)$. Ποια παράσταση ταιριάζει στο πρόβλημα και πόσο πληρώνει τελικά ο Άρης;

\ex \textbf{Τρεις απαντήσεις.} Τρεις φίλοι υπολογίζουν την παράσταση $4\cdot(7+7\cdot 9)+20$. Ο Άρης βρίσκει 335, η Ζωή 300 και ο Ηλίας 524. Ποιος έχει δίκιο; Βρες πού έκανε λάθος καθένας από τους άλλους δύο.

\zoni{Το σκέφτομαι}{μια ερώτηση που πάει πιο μακριά}

\ex Πόσο κάνει το $2^3\cdot 2^2$; Γράψε το ως γινόμενο παραγόντων (όλα τα δυάρια) και μέτρα. Ανακαλύπτεις έναν γρήγορο κανόνα για τον πολλαπλασιασμό δυνάμεων με την ίδια βάση;

\ex Συγκρίνε τα $2^5$ και $5^2$, και τα $2^4$ και $4^2$. Μετά βρες έναν αριθμό $\alpha$ διαφορετικό από το 2 για τον οποίο το $\alpha^2$ είναι ίσο με το $2\alpha$.

\ex Τοποθέτησε παρενθέσεις στην παράσταση $2+3\cdot 4+5$ με όσους τρόπους μπορείς και βρες όλες τις διαφορετικές τιμές που μπορεί να πάρει.

\ex Βάλε τα σύμβολα $+,-,\cdot$ (ή παρενθέσεις) ανάμεσα στα ψηφία $1\ \ 2\ \ 3\ \ 4$ ώστε το αποτέλεσμα να είναι 10. Βρες τουλάχιστον δύο τρόπους.

\ex Ο Ηλίας λέει: «Το $10^0$ πρέπει να είναι 0, γιατί δεν πολλαπλασιάζεις το 10 καμία φορά». Η Ζωή λέει: «Μάλλον είναι 1». Δες το μοτίβο $10^3=1.000$, $10^2=100$, $10^1=10$ και σκέψου τι πρέπει να συνεχίσει. Ποιος νομίζεις ότι έχει δίκιο; (Θα το μάθεις στο σχολείο σε λίγους μήνες.)

\zoni{Για το σπίτι}{εξάσκηση και ιστορία των αριθμών}

\ex Βρες τα τετράγωνα των 10, 20, 30, 40, 50, 60, 70, 80, 90 και τους κύβους των 10, 20, 30, 40, 50. Τι παρατηρείς στα μηδενικά;

\ex Υπολόγισε τις δυνάμεις $2^1, 2^2, 2^3, \dots, 2^{10}$ με τη σειρά, κάθε μία από την προηγούμενη.

\ex Γράψε πιο σύντομα: α) $\alpha+\alpha+\alpha$ \quad β) $\alpha\cdot\alpha\cdot\alpha$ \quad γ) $x+x+x+x$ \quad δ) $x\cdot x\cdot x\cdot x$ \quad ε) $2\cdot 2\cdot\alpha\cdot\alpha\cdot\alpha$.

\ex Βρες τις τιμές των παραστάσεων $(6+5)^2$ και $6^2+5^2$, και μετά των $(3+6)^2$ και $3^2+6^2$. Τι παρατηρείς;

\ex \textbf{Ρωμαϊκοί αριθμοί.} α) Γράψε με ρωμαϊκά τους 4, 9, 14, 40, 49, 90 και 1.994. β) Διάβασε τους XXXVII, CDXLIV και MMXXVI.

\ex \textbf{Αλφαβητικοί αριθμοί.} α) Γράψε με το αλφαβητικό σύστημα των Ιώνων τους 12, 25 και 831. β) Διάβασε τους ρκγ, σμ και τξε.

\ex \textbf{Δυαδικοί αριθμοί.} α) Γράψε στο δυαδικό σύστημα τους 11, 12 και 16. β) Διάβασε τους δυαδικούς 111, 1101 και 10101.

\ex \textbf{Οι αριθμοί των Πυθαγορείων.} Συνέχισε την ακολουθία των τριγώνων αριθμών $1,3,6,10$ για τρεις ακόμα όρους. Πόσο είναι ο δέκατος τρίγωνος αριθμός; Πόσο είναι ο δωδέκατος τετράγωνος;

\ex \textbf{Μαγικό τετράγωνο.} Βάλε τους αριθμούς από το 1 ως το 9, από μία φορά τον καθένα, ώστε κάθε γραμμή, κάθε στήλη και κάθε διαγώνιος να έχει άθροισμα 15. Έχουν μπει ήδη μερικοί:
\begin{center}
\renewcommand{\arraystretch}{1.6}\setlength{\tabcolsep}{12pt}
\begin{tabular}{|c|c|c|}\hline 2 & & 6\\\hline & 5 & \\\hline 4 & & 8\\\hline\end{tabular}
\end{center}

\paragrafos{ΕΚΠ και ΜΚΔ: πότε ξανασυναντιόμαστε, πόσο μεγάλα κομμάτια}{Α.1.4, «ΕΚΠ και ΜΚΔ φυσικών αριθμών»}

\begin{anagke}
Η γιαγιά Μυρτώ ποτίζει τα φυτά της αυλής κάθε 3 ημέρες και ο παππούς Γιάννης κλαδεύει τη λεμονιά κάθε 4 ημέρες. Σήμερα έκαναν και οι δύο τη δουλειά τους. «Πότε θα τα ξανακάνουμε και οι δύο την ίδια μέρα;» ρωτά η Ζωή, που έχει έρθει για επίσκεψη. Ο Ηλίας γράφει στο τετράδιο τις ημέρες της γιαγιάς: 3, 6, 9, 12, 15, 18, και του παππού: 4, 8, 12, 16, 20. «Στις 12!» φωνάζει.

Το απόγευμα η θεία Μαρία φέρνει 48 μανταρίνια και 60 πορτοκάλια και θέλει να τα βάλει σε ίσα καλάθια, με τον ίδιο αριθμό μανταρινιών και πορτοκαλιών σε κάθε καλάθι, χωρίς να περισσέψει τίποτα. «Πόσα το πολύ καλάθια μπορώ να φτιάξω;» ρωτά. Αυτό δεν είναι «πότε ξανασυναντιόμαστε»: είναι «σε πόσα ίσα κομμάτια χωρίζονται και τα δύο».

Δύο διαφορετικά προβλήματα: το ένα ψάχνει έναν αριθμό που τον φτάνουν και οι δύο (πολλαπλάσιο), το άλλο έναν αριθμό που χωράει σε όλους (διαιρέτη). Σε αυτή την παράγραφο μαθαίνουμε το ελάχιστο κοινό πολλαπλάσιο (ΕΚΠ) και τον μέγιστο κοινό διαιρέτη (ΜΚΔ), και πώς τα βρίσκουμε με την ανάλυση σε γινόμενο πρώτων παραγόντων και με τον αλγόριθμο του Ευκλείδη.
\end{anagke}

\begin{lexi}{πολλαπλάσιο, κοινό, ελάχιστο, μέγιστο}{πολλαπλάσια, ΕΚΠ, ΜΚΔ}
Το \lexiword{πολλαπλάσιο} ενός αριθμού είναι αυτός που προκύπτει όταν τον πάρουμε πολλές φορές, δηλαδή τον πολλαπλασιάσουμε. \lexiword{Κοινό} είναι αυτό που ανήκει και στους δύο. Με το ΕΚΠ ψάχνουμε το \lexiword{ελάχιστο}, το πιο μικρό, κοινό πολλαπλάσιο, και με το ΜΚΔ τον \lexiword{μέγιστο}, τον πιο μεγάλο, κοινό διαιρέτη. Το $0$ είναι πολλαπλάσιο κάθε αριθμού, γι' αυτό στο ΕΚΠ ζητάμε το μικρότερο \emph{μη μηδενικό} κοινό πολλαπλάσιο.
\end{lexi}

\begin{metafrasi}
Οι αριθμοί που προκύπτουν όταν πολλαπλασιάσω το 4. & $0,4,8,12,16,\dots$ (πολλαπλάσια του $4$)\\
Το πρώτο κοινό πολλαπλάσιο των 4 και 6. & $\EKP(4,6)=12$\\
Το μεγαλύτερο κοινό τους κομμάτι των 12 και 18. & $\MKD(12,18)=6$\\
Δύο αριθμοί δεν έχουν κοινό διαιρέτη εκτός από το 1. & $\MKD(\alpha,\beta)=1$, πρώτοι μεταξύ τους\\
\end{metafrasi}

\tmima{Η θεωρία}

\tmimasmall{Πολλαπλάσια και ΕΚΠ}

\begin{orismos}{πολλαπλάσια}
\textbf{Πολλαπλάσια} ενός φυσικού $\alpha$ είναι οι αριθμοί που προκύπτουν όταν τον πολλαπλασιάσουμε με όλους τους φυσικούς: $0,\ \alpha,\ 2\alpha,\ 3\alpha,\ 4\alpha,\dots$ Για παράδειγμα τα πολλαπλάσια του $9$ είναι τα $0,9,18,27,36,45,\dots$
\end{orismos}

\begin{orismos}{ελάχιστο κοινό πολλαπλάσιο}
Το \textbf{ελάχιστο κοινό πολλαπλάσιο} (ΕΚΠ) δύο ή περισσότερων φυσικών αριθμών που δεν είναι μηδέν είναι το μικρότερο (μη μηδενικό) από τα κοινά πολλαπλάσια των αριθμών αυτών.
\end{orismos}

Παράδειγμα: τα πολλαπλάσια του $4$ είναι $0,4,8,\mathbf{12},16,20,\mathbf{24},28,32,\mathbf{36},\dots$ και του $6$ είναι $0,6,\mathbf{12},18,\mathbf{24},30,\mathbf{36},42,\dots$ Τα κοινά είναι τα $12,24,36,\dots$ (άπειρα) και το μικρότερο είναι το $12$, άρα $\EKP(4,6)=12$.

\tmimasmall{Διαιρέτες και ΜΚΔ}

Από την προηγούμενη παράγραφο ξέρουμε ότι \textbf{διαιρέτες} ενός αριθμού είναι όσοι τον διαιρούν τέλεια. Ισχύουν τα εξής:
\begin{itemize}[leftmargin=1.2em,itemsep=0.2ex,topsep=2pt]
\item Κάθε φυσικός (εκτός από το $0$) διαιρεί τα πολλαπλάσιά του.
\item Κάθε φυσικός που διαιρείται από έναν άλλον είναι πολλαπλάσιο αυτού του αριθμού.
\item Αν ένας φυσικός διαιρεί έναν άλλον, θα διαιρεί και τα πολλαπλάσια του άλλου.
\end{itemize}

\begin{orismos}{μέγιστος κοινός διαιρέτης}
Ο \textbf{μέγιστος κοινός διαιρέτης} (ΜΚΔ) δύο ή περισσότερων φυσικών αριθμών είναι ο μεγαλύτερος από τους κοινούς διαιρέτες τους.
\end{orismos}

Παράδειγμα: οι διαιρέτες του $12$ είναι οι $1,2,3,4,6,12$ και του $18$ οι $1,2,3,6,9,18$. Οι κοινοί διαιρέτες είναι οι $1,2,3,6$ και ο μεγαλύτερος είναι ο $6$, άρα $\MKD(12,18)=6$.

Αν $\MKD(\alpha,\beta)=1$, οι αριθμοί $\alpha$ και $\beta$ λέγονται \textbf{πρώτοι μεταξύ τους}. Για παράδειγμα οι $6$ και $7$ έχουν κοινό διαιρέτη μόνο το $1$, άρα είναι πρώτοι μεταξύ τους (αν και ο $6$ δεν είναι πρώτος αριθμός).

\begin{prosoxi}
Το «πρώτοι μεταξύ τους» δεν σημαίνει ότι οι δύο αριθμοί είναι πρώτοι. Οι $8$ και $15$ δεν είναι πρώτοι αριθμοί, αλλά είναι πρώτοι μεταξύ τους.
\end{prosoxi}

\tmimasmall{ΕΚΠ και ΜΚΔ με ανάλυση σε πρώτους παράγοντες}

Η λίστα με τα πολλαπλάσια ή τους διαιρέτες είναι κουραστική για μεγάλους αριθμούς. Αναλύουμε πρώτα όλους τους αριθμούς σε γινόμενο πρώτων παραγόντων.

\begin{kanonas}{ΕΚΠ και ΜΚΔ με ανάλυση}
Το \textbf{ΕΚΠ} είναι το γινόμενο όλων των πρώτων παραγόντων (κοινών και μη κοινών) με τον \textbf{μεγαλύτερο} εκθέτη τους.\\
Ο \textbf{ΜΚΔ} είναι το γινόμενο μόνο των \textbf{κοινών} πρώτων παραγόντων με τον \textbf{μικρότερο} εκθέτη τους.
\end{kanonas}

\tmimasmall{Ο αλγόριθμος του Ευκλείδη για το ΜΚΔ}

Ο Ευκλείδης έδωσε μια μέθοδο που βρίσκει τον ΜΚΔ χωρίς καμία ανάλυση. Τα βήματα, για τους αριθμούς $125,170,280$:

\begin{center}
\renewcommand{\arraystretch}{1.3}\setlength{\tabcolsep}{9pt}
\begin{tabular}{ccc}
$125$ & $170$ & $280$\\
$125$ & $45$ & $30$\\
$5$ & $15$ & $30$\\
$5$ & $0$ & $0$
\end{tabular}
\end{center}

Γράφουμε τους αριθμούς και βρίσκουμε τον μικρότερο. Στην επόμενη γραμμή γράφουμε \emph{τον μικρότερο αριθμό ως έχει} και κάτω από κάθε άλλον το \emph{υπόλοιπο της διαίρεσής του με τον μικρότερο} ($170=125\cdot 1+45$, $280=125\cdot 2+30$). Επαναλαμβάνουμε με τον νέο μικρότερο μη μηδενικό αριθμό. Σταματάμε όταν όλοι οι άλλοι γίνουν $0$. Ο ΜΚΔ είναι ο μη μηδενικός αριθμός της τελευταίας γραμμής: $\MKD(125,170,280)=5$.

\tmima{Λυμένα παραδείγματα}

\paradeigma{Το λεωφορείο και το τρένο}
\tixero Από το χωριό φεύγει λεωφορείο κάθε $12$ λεπτά και τρένο κάθε $15$ λεπτά. Έφυγαν μαζί στις $8{:}00$.\par
\tithelo Πότε θα ξαναφύγουν μαζί;\par
\possindeo Ζητάω το μικρότερο κοινό πολλαπλάσιο των $12$ και $15$. Πολλαπλάσια του $12$: $12,24,36,48,60,\dots$ Πολλαπλάσια του $15$: $15,30,45,60,\dots$ Το πρώτο κοινό είναι το $60$.\par
\apantisi $\EKP(12,15)=60$, άρα μετά από $60$ λεπτά, στις $9{:}00$.\par
\elegxos Το $60=12\cdot 5=15\cdot 4$.

\paradeigma{Οι χάντρες της Αγγελικής}
\tixero Η Αγγελική έχει $36$ κόκκινες, $48$ κίτρινες και $60$ μπλε χάντρες. Φτιάχνει όσο το δυνατόν περισσότερα ίδια κολιέ, με τον ίδιο αριθμό χαντρών κάθε χρώματος, χωρίς να περισσέψει καμία.\par
\tithelo Πόσα κολιέ φτιάχνει και τι έχει το καθένα;\par
\possindeo Ο αριθμός των κολιέ πρέπει να διαιρεί και τα τρία πλήθη, και ζητάω τον μεγαλύτερο, άρα τον ΜΚΔ. Αναλύω: $36=2^2\cdot 3^2$, $48=2^4\cdot 3$, $60=2^2\cdot 3\cdot 5$. Κοινοί παράγοντες: $2$ και $3$ με μικρότερους εκθέτες $2$ και $1$. Άρα $\MKD=2^2\cdot 3=12$. Κάθε κολιέ έχει $36:12=3$ κόκκινες, $48:12=4$ κίτρινες και $60:12=5$ μπλε.\par
\apantisi $12$ κολιέ, με $3$ κόκκινες, $4$ κίτρινες και $5$ μπλε χάντρες.\par
\elegxos $12\cdot 3=36$, $12\cdot 4=48$, $12\cdot 5=60$.

\paradeigma{ΕΚΠ και ΜΚΔ με ανάλυση, και με Ευκλείδη}
\tixero Οι αριθμοί $84$, $90$ και $126$.\par
\tithelo Το ΕΚΠ και ο ΜΚΔ τους.\par
\possindeo Αναλύω: $84=2^2\cdot 3\cdot 7$, $90=2\cdot 3^2\cdot 5$, $126=2\cdot 3^2\cdot 7$. Όλοι οι πρώτοι παράγοντες είναι οι $2,3,5,7$, με μεγαλύτερους εκθέτες $2,2,1,1$: $\EKP=2^2\cdot 3^2\cdot 5\cdot 7=1.260$. Κοινοί παράγοντες είναι μόνο οι $2$ και $3$, με μικρότερους εκθέτες $1,1$: $\MKD=2\cdot 3=6$. Με τον αλγόριθμο του Ευκλείδη: $(84,90,126)\to(84,6,42)\to(0,6,0)$, άρα ΜΚΔ $=6$.\par
\apantisi $\EKP=1.260$ και $\MKD=6$.\par
\elegxos Το $1.260:84=15$, $1.260:90=14$, $1.260:126=10$ και ο $6$ διαιρεί τους $84,90,126$.

\tmima{Ασκήσεις}

\zoni{Το κατάλαβα}{ορισμοί και υπολογισμοί}

\ex Συμπλήρωσε: α) Αν $\MKD(\alpha,\beta)=1$, οι $\alpha$ και $\beta$ λέγονται \dots μεταξύ τους. β) Κάθε φυσικός \dots τα πολλαπλάσιά του. γ) Κάθε φυσικός που διαιρείται από έναν άλλον είναι \dots αυτού του αριθμού. δ) Ο ΜΚΔ δύο αριθμών είναι ο \dots από τους κοινούς διαιρέτες τους.

\ex Γράψε τα $8$ πρώτα πολλαπλάσια των $7$ και $12$ (ξεκινώντας από το $0$).

\ex Γράψε τα $8$ πρώτα μη μηδενικά πολλαπλάσια των $6$ και $9$. Ποια είναι κοινά; Ποιο είναι το $\EKP(6,9)$;

\ex Γράψε όλους τους διαιρέτες των $18$, $30$ και $36$. Βρες τους κοινούς διαιρέτες και τον $\MKD(18,30,36)$.

\ex Είναι πρώτοι μεταξύ τους οι αριθμοί; α) $8$ και $15$ \quad β) $14$ και $21$ \quad γ) $22$ και $45$ \quad δ) $9$ και $28$ \quad ε) $12$ και $35$.

\ex Συμπλήρωσε τον πίνακα με τον ΜΚΔ και το ΕΚΠ κάθε ζεύγους: $(6,8)$, $(10,15)$, $(4,9)$, $(14,21)$, $(16,24)$.

\ex Ανάλυσε τους $36$ και $54$ σε γινόμενο πρώτων παραγόντων και βρες τον ΜΚΔ και το ΕΚΠ τους.

\zoni{Το χρησιμοποιώ}{ΕΚΠ ή ΜΚΔ;}

\ex \textbf{Οι προπονήσεις.} Ο Νίκος προπονείται κάθε $3$ ημέρες, ο Βασίλης κάθε $5$ και η Πολυξένη κάθε $6$. Σήμερα προπονήθηκαν όλοι μαζί. Μετά από πόσες ημέρες θα ξανασυναντηθούν;

\ex \textbf{Τα καλάθια.} Η θεία Μαρία έχει $54$ μήλα και $72$ αχλάδια και τα βάζει σε ίσα καλάθια, χωρίς να περισσέψει τίποτα και με τον ίδιο αριθμό μήλων και αχλαδιών σε κάθε καλάθι. Πόσα το πολύ καλάθια μπορεί να φτιάξει και τι έχει το καθένα;

\ex \textbf{Τα πλακάκια.} Το δωμάτιο έχει διαστάσεις $240\,\mathrm{cm}\times 180\,\mathrm{cm}$ και θέλουμε να το στρώσουμε με τετράγωνα πλακάκια, ίδια, χωρίς κοψίματα. Ποια είναι η μεγαλύτερη δυνατή πλευρά του πλακιδίου και πόσα πλακάκια χρειάζονται;

\ex \textbf{Τα φωτάκια.} Στη γιορτή του χωριού ένα κόκκινο φωτάκι ανάβει κάθε $12$ δευτερόλεπτα και ένα πράσινο κάθε $18$. Άναψαν μαζί. Μετά από πόσα δευτερόλεπτα θα ανάψουν ξανά μαζί;

\ex \textbf{Οι σελιδοδείκτες.} Η Φανή έχει $40$ σελιδοδείκτες και $56$ αυτοκόλλητα. Φτιάχνει ίδια πακέτα, με τον ίδιο αριθμό από κάθε είδος, χωρίς να περισσέψει τίποτα. Πόσα το πολύ πακέτα φτιάχνει και τι έχει κάθε πακέτο;

\ex \textbf{Ο φούρνος.} Ένας φούρνος έχει $60$ κουλούρια, $84$ κρουασάν και $108$ τσουρέκια και τα βάζει σε ίδια κουτιά, με τον ίδιο αριθμό από κάθε είδος και χωρίς να περισσέψει τίποτα. Πόσα το πολύ κουτιά γεμίζει και τι έχει το καθένα;

\zoni{Το σκέφτομαι}{ερωτήσεις με λίγη περισσότερη σκέψη}

\ex \textbf{Μια ανακάλυψη.} Για τα ζεύγη $(4,6)$, $(12,18)$ και $(10,15)$ βρες τον ΜΚΔ και το ΕΚΠ. Μετά υπολόγισε το γινόμενο $\MKD\cdot\EKP$ και το γινόμενο των δύο αριθμών κάθε ζεύγους. Τι παρατηρείς;

\ex Βρες δύο διαφορετικούς φυσικούς που έχουν $\MKD=6$ και $\EKP=60$.

\ex α) Βρες το $\EKP(8,15)$ και το $\EKP(9,10)$. Πώς σχετίζονται με το γινόμενο των αριθμών, όταν οι αριθμοί είναι πρώτοι μεταξύ τους; β) Βρες τον $\MKD$ και το $\EKP$ των $6$ και $30$ (ο ένας διαιρεί τον άλλον). Τι παρατηρείς;

\zoni{Για το σπίτι}{εξάσκηση}

\ex Βρες με ανάλυση τον ΜΚΔ και το ΕΚΠ των $325$ και $338$ (από την παράγραφο 1.3) και των $24$, $42$, $54$.

\ex Βρες με τον αλγόριθμο του Ευκλείδη τον ΜΚΔ των $48$, $80$, $120$. Μετά εξέτασε αν οι $1.000$ και $1.303$ είναι πρώτοι μεταξύ τους.

\paragrafos{Χαρακτήρες διαιρετότητας, πρώτοι αριθμοί, ΜΚΔ και ΕΚΠ}{Α.1.5, «Χαρακτήρες διαιρετότητας – ΜΚΔ – ΕΚΠ – Ανάλυση αριθμού σε γινόμενο πρώτων παραγόντων»}
\ypotmima{Α}{Χαρακτήρες διαιρετότητας: ξέρω χωρίς να διαιρέσω}


\begin{anagke}
Το σχολικό παζάρι πήγε καλύτερα απ' όσο περίμενε κανείς. Η Φανή, που κρατάει το ταμείο, μετράει τα χρήματα: 2.358 €. Θέλουν να τα μοιράσουν ίσα σε τρεις ομάδες που θα φτιάξουν από ένα σχολικό έργο. Μετά η Ρήνη λέει: «Και αν οι ομάδες ήταν εννιά;» Ο Νίκος: «Και αν ήταν πέντε;» Ο Στέλιος ξεκινάει να κάνει τις διαιρέσεις με τη σειρά, μία μία. Τρεις διαιρέσεις με τετραψήφιο αριθμό είναι πολλή δουλειά.

Η Πολυξένη κοιτάζει τον αριθμό 2.358 και λέει: «Για το πέντε δεν χρειάζεται καν να κάνεις διαίρεση. Τελειώνει σε οκτώ. Δεν γίνεται». «Πώς το ξέρεις;» ρωτάει ο Ηλίας. «Το έχω δει σε όλους τους αριθμούς που διαιρούνται με το πέντε: τελειώνουν σε 0 ή σε 5». Η Ζωή ρωτάει: «Και για το τρία και για το εννιά υπάρχει κάτι τέτοιο;»

Υπάρχει. Και είναι από τα πιο χρήσιμα κόλπα στο γυμνάσιο: να καταλαβαίνεις αν ένας αριθμός διαιρείται με έναν άλλον κοιτάζοντας απλώς τα ψηφία του.
\end{anagke}

\begin{lexi}{κριτήριο, χαρακτήρας διαιρετότητας}{κριτήριο, χαρακτήρας}
Το \lexiword{κριτήριο} έρχεται από το «κρίνω», δηλαδή αποφασίζω. Κριτήριο είναι ο κανόνας με τον οποίο κρίνουμε κάτι. Ένας κριτής σε αγώνα έχει κριτήρια για να βγάλει τον νικητή. Ένα \lexiword{κριτήριο διαιρετότητας} μας βοηθάει να κρίνουμε αν ένας αριθμός διαιρείται με κάποιον άλλον, χωρίς να κάνουμε τη διαίρεση.

Το σχολικό βιβλίο τα λέει και \lexiword{χαρακτήρες διαιρετότητας}. Ο χαρακτήρας είναι το «σημάδι» που χαρακτηρίζει κάτι, όπως ο χαρακτήρας ενός ανθρώπου. Εδώ το σημάδι είναι στα ψηφία του αριθμού.
\end{lexi}

\begin{metafrasi}
Ο αριθμός είναι άρτιος. & Λήγει σε $0,\,2,4,\,6$ ή $8$ (διαιρείται με το 2)\\
Ο αριθμός διαιρείται με το 5. & Λήγει σε $0$ ή $5$\\
Ο αριθμός διαιρείται με το 10. & Λήγει σε $0$\\
Ο αριθμός διαιρείται με το 3. & Το άθροισμα των ψηφίων του διαιρείται με το 3\\
Ο αριθμός διαιρείται με το 9. & Το άθροισμα των ψηφίων του διαιρείται με το 9\\
Ο αριθμός διαιρείται με το 4. & Ο αριθμός από τα δύο τελευταία ψηφία διαιρείται με το 4\\
Ο αριθμός διαιρείται με το 25. & Λήγει σε $00$, $25$, $50$ ή $75$\\
Ο αριθμός διαιρείται με το 6. & Διαιρείται και με το 2 και με το 3\\
\end{metafrasi}

\tmima{Η θεωρία}

\begin{kanonas}{κριτήρια διαιρετότητας}
\begin{itemize}[leftmargin=1.2em,itemsep=0.3ex,topsep=2pt]
\item \textbf{Με το 2}: το τελευταίο ψηφίο είναι 0, 2, 4, 6 ή 8.
\item \textbf{Με το 5}: το τελευταίο ψηφίο είναι 0 ή 5.
\item \textbf{Με το 10}: το τελευταίο ψηφίο είναι 0. (Με το 100: τα δύο τελευταία ψηφία είναι 00, με το 1.000: τα τρία τελευταία, κ.ο.κ.)
\item \textbf{Με το 4}: ο αριθμός που σχηματίζουν τα δύο τελευταία ψηφία διαιρείται με το 4 (ή είναι 00).
\item \textbf{Με το 25}: τα δύο τελευταία ψηφία είναι 00, 25, 50 ή 75.
\item \textbf{Με το 3}: το άθροισμα των ψηφίων διαιρείται με το 3.
\item \textbf{Με το 9}: το άθροισμα των ψηφίων διαιρείται με το 9.
\end{itemize}
\end{kanonas}

Από τα κριτήρια αυτά φτιάχνουμε και άλλα. Για να διαιρείται ένας αριθμός με το 6 πρέπει να διαιρείται και με το 2 και με το 3. Για το 15, και με το 3 και με το 5. Για το 12, και με το 3 και με το 4.

\tmimasmall{Γιατί δουλεύουν}

Κάθε αριθμός ξεκινάει με πολλαπλάσια του 10, του 100, του 1.000. Για παράδειγμα $3.725=37\cdot 100+25$. Ο 100 διαιρείται με το 4 και με το 25, άρα το $37\cdot 100$ διαιρείται και με τους δύο, και το μόνο που μένει να ελέγξουμε είναι τα δύο τελευταία ψηφία, το 25. Το ίδιο και για τα 2 και 5, με το 10.

Για το 3 και το 9 δουλεύει ένα κόλπο. Οι αριθμοί 9, 99, 999, \dots\ διαιρούνται με το 9 (και με το 3). Γράφω $2.457=2\cdot 1.000+4\cdot 100+5\cdot 10+7$ και αντικαθιστώ το $1.000=999+1$, $100=99+1$, $10=9+1$:
$$2.457=2\cdot 999+4\cdot 99+5\cdot 9+(2+4+5+7).$$
Τα τρία πρώτα κομμάτια διαιρούνται με το 9. Άρα όλος ο αριθμός διαιρείται με το 9 ακριβώς όταν διαιρείται το $2+4+5+7=18$. Και διαιρείται.

\begin{prosoxi}
\begin{itemize}[leftmargin=1.2em,itemsep=0.3ex,topsep=2pt]
\item Αν ένας αριθμός τελειώνει σε 4, δεν σημαίνει ότι διαιρείται με το 4. Το 14 δεν διαιρείται. Για το 4 κοιτάμε τα δύο τελευταία ψηφία μαζί.
\item Για το 3 και το 9 μετράνε \textbf{όλα} τα ψηφία, όχι μόνο το τελευταίο.
\item Για να διαιρείται ένας αριθμός με το 24 δεν φτάνει να διαιρείται με το 4 και με το 6. Το 12 διαιρείται και με τα δύο, αλλά όχι με το 24. Μπορούμε να «σπάσουμε» έναν αριθμό σε δύο μόνο όταν οι δύο αριθμοί δεν έχουν κοινό παράγοντα (όπως το 3 και το 4, ή το 2 και το 5).
\item Το 0 διαιρείται με κάθε αριθμό (εκτός από το 0 ως διαιρέτη), γιατί $0=\delta\cdot 0$.
\end{itemize}
\end{prosoxi}

\tmima{Λυμένα παραδείγματα}

\paradeigma{Ο πίνακας των κριτηρίων}
\tixero Οι αριθμοί 4.860, 975, 1.236 και 7.500.\par
\tithelo Με ποιους από τους 2, 3, 4, 5, 9, 10, 25 και 100 διαιρείται ο καθένας.\par
\possindeo Εφαρμόζω κάθε κριτήριο. Στα 3 και 9 υπολογίζω το άθροισμα ψηφίων. Στα 4 και 25 κοιτάζω τα δύο τελευταία ψηφία.

\begin{center}
\small\renewcommand{\arraystretch}{1.3}\setlength{\tabcolsep}{4pt}
\begin{tabular}{@{}lcccccccc@{}}
\toprule
 & 2 & 3 & 4 & 5 & 9 & 10 & 25 & 100\\
\midrule
4.860 \ ($4+8+6+0=18$) & ναι & ναι & ναι & ναι & ναι & ναι & όχι & όχι\\
975 \ ($9+7+5=21$) & όχι & ναι & όχι & ναι & όχι & όχι & ναι & όχι\\
1.236 \ ($1+2+3+6=12$) & ναι & ναι & ναι & όχι & όχι & όχι & όχι & όχι\\
7.500 \ ($7+5=12$) & ναι & ναι & ναι & ναι & όχι & ναι & ναι & ναι\\
\bottomrule
\end{tabular}
\end{center}

\apantisi Δες τον πίνακα.\par
\elegxos Το 7.500 διαιρείται με το 100, άρα και με το 10, το 4, το 25, το 5 και το 2: κάθε διαιρέτης του 100 μπαίνει αυτόματα. Δοκιμή με το 4: $7.500:4=1.875$. Σωστά. Για το 1.236 με το 4: $36:4=9$, και πράγματι $1.236:4=309$.

\paradeigma{Το ψηφίο που λείπει}
\tixero Ο αριθμός $3.8\square 1$, όπου το $\square$ είναι ένα ψηφίο.\par
\tithelo Ποιο ψηφίο πρέπει να μπει ώστε ο αριθμός να διαιρείται με το 9.\par
\possindeo Το άθροισμα των γνωστών ψηφίων είναι $3+8+1=12$. Πρέπει $12+\square$ να διαιρείται με το 9. Το επόμενο πολλαπλάσιο του 9 μετά το 12 είναι το 18, άρα $\square=6$. (Το επόμενο, 27, θα απαιτούσε ψηφίο 15, που δεν υπάρχει.)\par
\apantisi $3.861$.\par
\elegxos $3.861:9=429$, γιατί $9\cdot 429=3.861$. Σωστά.

\paradeigma{Η μοιρασιά του παζαριού}
\tixero Η Φανή έχει 2.358 €.\par
\tithelo Σε ποιες από τις 3, 9, 5, 4 και 6 ισόποσες ομάδες μοιράζονται χωρίς να περισσεύει τίποτα.\par
\possindeo Άθροισμα ψηφίων: $2+3+5+8=18$. Διαιρείται με το 3 και με το 9. Τελευταίο ψηφίο 8, άρα όχι με το 5. Τελευταία δύο ψηφία 58: το $58:4$ δεν είναι ακέραιο, άρα όχι με το 4. Είναι άρτιος και διαιρείται με το 3, άρα διαιρείται με το 6.\par
\apantisi Σε 3, 6 και 9 ομάδες.\par
\elegxos $2.358:3=786$, $2.358:6=393$, $2.358:9=262$. Και οι τρεις είναι τέλειες διαιρέσεις.

\tmima{Ασκήσεις}

\zoni{Το κατάλαβα}{εφαρμογή του κανόνα}

\ex Ποιοι από τους αριθμούς 2.480, 3.115, 6.002, 1.050 και 819 διαιρούνται με το 2; Ποιοι με το 5; Ποιοι με το 10;

\ex Βρες ποιοι από τους 252, 3.614, 8.001, 55.551 και 1.370 διαιρούνται με το 3. Έλεγξε μία από τις περιπτώσεις με διαίρεση.

\ex Βρες ποιοι από τους 738, 5.427, 92.808, 6.661 και 14.445 διαιρούνται με το 9.

\ex Βρες ποιοι από τους 524, 1.018, 6.300, 3.772 και 90 διαιρούνται με το 4.

\ex Βρες ποιοι από τους 450, 2.125, 3.300, 1.080 και 7.675 διαιρούνται με το 25.

\ex Συμπλήρωσε έναν πίνακα με τους αριθμούς 360, 1.125, 5.724 και 40.500 και τους διαιρέτες 2, 3, 4, 5, 9, 10, 25. Για κάθε αριθμό γράψε ναι ή όχι.

\ex Ποιοι από τους 78, 250, 3.402, 1.113 και 9.996 διαιρούνται με το 6;

\ex Βρες το ψηφίο (ή τα ψηφία) που πρέπει να μπουν στη θέση του $\square$: α) $6.3\square 8$ να διαιρείται με το 9 \quad β) $4\square 2$ να διαιρείται με το 3 \quad γ) $2\square 6$ να διαιρείται με το 4 \quad δ) $29\square$ να διαιρείται και με το 2 και με το 5.

\ex Σωστό ή λάθος; α) Κάθε αριθμός που τελειώνει σε 6 διαιρείται με το 6. \quad β) Αν ένας αριθμός διαιρείται με το 9, διαιρείται και με το 3. \quad γ) Το 0 διαιρείται με το 13. \quad δ) Αν ένας αριθμός διαιρείται με το 6 και με το 8, διαιρείται και με το 48. \quad ε) Αν διαιρείται με το 3 και με το 5, διαιρείται με το 15. \quad στ) Αν διαιρείται με το 4 και με το 25, διαιρείται με το 100.

\zoni{Το χρησιμοποιώ}{προβλήματα από τη ζωή}

\ex \textbf{Τα έσοδα της λαχειοφόρου.} Η λαχειοφόρος του σχολείου έφερε 4.752 €. Σε ποιες από τις 3, 4, 5, 6 και 9 ομάδες μοιράζονται ακριβώς; Βρες πόσα παίρνει η κάθε ομάδα, όπου γίνεται.

\ex \textbf{Τα μπλουζάκια.} Ο προπονητής του Νίκου έχει 150 μπλουζάκια και θέλει να τα βάλει σε ίσες στοίβες. Με ποιους από τους αριθμούς 2, 3, 4, 5, 6 και 10 γίνεται, χωρίς να περισσεύει κανένα; Χρησιμοποίησε κριτήρια.

\ex \textbf{Ο κωδικός της Πολυξένης.} Ο νέος κωδικός της Πολυξένης είναι τετραψήφιος. Ξεκινάει με 5, το δεύτερο ψηφίο είναι το 4, τελειώνει σε 0, διαιρείται με το 9 και το τρίτο ψηφίο είναι μεγαλύτερο από το δεύτερο. Ποιος είναι ο κωδικός;

\ex \textbf{Οι θέσεις του γηπέδου.} Οι θέσεις ενός μικρού γηπέδου είναι αριθμημένες από το 1 ως το 90. Οι θέσεις που διαιρούνται με το 5 έχουν σκιά. Πόσες είναι; Πόσες διαιρούνται με το 10; Πόσες διαιρούνται και με το 3 και με το 5;

\ex \textbf{Τα δίσεκτα έτη.} Ένα έτος είναι δίσεκτο αν διαιρείται με το 4. Εξαίρεση οι αιώνες (1.900, 2.000, \dots), που είναι δίσεκτοι μόνο όταν διαιρούνται με το 400. Ποια από τα έτη 1.900, 2.000, 2.024, 2.030 και 2.100 είναι δίσεκτα;

\zoni{Το σκέφτομαι}{μια ερώτηση που πάει πιο μακριά}

\ex Γιατί ένας αριθμός που τελειώνει σε 0 διαιρείται και με το 2 και με το 5; (Βοήθεια: $10=2\cdot 5$.)

\ex Γράψε το 2.763 ως $2\cdot 999+7\cdot 99+6\cdot 9+(2+7+6+3)$ και εξήγησε με τη βοήθεια αυτής της γραφής αν διαιρείται με το 3 και αν διαιρείται με το 9.

\ex Ποιος είναι ο μικρότερος τετραψήφιος αριθμός που διαιρείται ταυτόχρονα με τα 4, 5 και 6;

\ex Οι αριθμοί 121, 132, 143 και 154 διαιρούνται όλοι με το 11. Τι παρατηρείς στα ψηφία τους; Δοκίμασε το ίδιο μοτίβο στο 165 και στο 175 και έλεγξε με διαίρεση.

\ex Ο Στέλιος ισχυρίζεται ότι ο 5.034 διαιρείται με το 4, γιατί τελειώνει σε 4. Ποιο είναι το λάθος του; Τι θα έπρεπε να είχε κάνει;

\ypotmima{Β}{Πρώτοι αριθμοί και ανάλυση: τα δομικά υλικά}

\begin{anagke}
Στο διάλειμμα η Πολυξένη, που κρατάει ένα τετράδιο με κωδικούς και μοτίβα, δείχνει στην παρέα κάτι που της εξήγησε η μαμά της, που δουλεύει στην ασφάλεια ενός δικτύου υπολογιστών. «Σκέφτομαι δύο πρώτους αριθμούς», λέει στη Ρήνη. «Το γινόμενό τους είναι 143. Ποιοι είναι;»

Η Ρήνη δοκιμάζει. Δεν διαιρείται με το 2, γιατί τελειώνει σε 3. Το άθροισμα ψηφίων $1+4+3=8$, άρα ούτε με το 3. Ούτε με το 5. Με το 7, $7\cdot 20=140$ και περισσεύουν 3, όχι. Με το 11: $11\cdot 13=143$. «Έντεκα και δεκατρία!» φωνάζει. «Μπράβο», λέει η Πολυξένη. «Τώρα σκέψου ότι ο αριθμός έχει 600 ψηφία. Ένας υπολογιστής πολλαπλασιάζει τους δύο πρώτους σε ένα κλάσμα του δευτερολέπτου. Για να τους βρει από το γινόμενο, θέλει απίστευτα πολύ χρόνο. Πάνω σε αυτό στηρίζονται οι κωδικοί των τραπεζών».

Η Ζωή ρωτάει κάτι πιο απλό. «Δηλαδή κάθε αριθμός είναι φτιαγμένος από πρώτους, όπως ο τοίχος από τούβλα;» Ναι. Και το να βρίσκεις αυτά τα τούβλα είναι δουλειά με μέθοδο.
\end{anagke}

\begin{lexi}{πρώτος, σύνθετος, παράγοντας}{πρώτος, σύνθετος}
Ο \lexiword{πρώτος} αριθμός δεν λέγεται έτσι επειδή έρχεται πρώτος στη σειρά, αλλά επειδή είναι «πρωταρχικός», πρωτογενής: δεν φτιάχνεται από μικρότερους αριθμούς με πολλαπλασιασμό. Είναι το δομικό υλικό.

Ο \lexiword{σύνθετος} αριθμός είναι αυτός που συντίθεται, δηλαδή φτιάχνεται, από άλλους: το 12 είναι $2\cdot 2\cdot 3$. Ο \lexiword{παράγοντας} είναι ο αριθμός που παίρνει μέρος σε έναν πολλαπλασιασμό και «παράγει» το γινόμενο.
\end{lexi}

\begin{metafrasi}
Το 7 έχει διαιρέτες μόνο το 1 και το 7. & Το 7 είναι πρώτος\\
Το 12 έχει διαιρέτες τους 1, 2, 3, 4, 6, 12. & Το 12 είναι σύνθετος\\
Το 60 σαν γινόμενο πρώτων. & $60=2\cdot 2\cdot 3\cdot 5=2^2\cdot 3\cdot 5$\\
Ο $\alpha$ διαιρεί τον $\beta$. & $\beta=\alpha\cdot\kappa$, με $\kappa$ φυσικό\\
Πολλαπλάσια του 4. & $0,\,4,\,8,\,12,\,16,\,20,\,\dots$\\
\end{metafrasi}

\tmima{Η θεωρία}

\tmimasmall{Πολλαπλάσια και διαιρέτες}

\textbf{Πολλαπλάσια} ενός φυσικού $\alpha$ είναι οι αριθμοί που βγαίνουν αν τον πολλαπλασιάσουμε με φυσικούς: $0,\ \alpha,\ 2\alpha,\ 3\alpha,\ 4\alpha,\ \dots$ \textbf{Διαιρέτες} του $\alpha$ είναι όλοι οι αριθμοί που τον διαιρούν ακριβώς. Κάθε αριθμός $\alpha$ έχει διαιρέτες τουλάχιστον τον 1 και τον εαυτό του. Αν ο $\delta$ διαιρεί έναν αριθμό, τότε ο αριθμός είναι πολλαπλάσιο του $\delta$. Είναι το ίδιο πράγμα, από δύο πλευρές.

\begin{orismos}{πρώτος και σύνθετος}
\textbf{Πρώτος} λέγεται ένας φυσικός αριθμός μεγαλύτερος από το 1 που έχει διαιρέτες μόνο τον 1 και τον εαυτό του. Ένας φυσικός μεγαλύτερος από το 1 που έχει και άλλους διαιρέτες λέγεται \textbf{σύνθετος}. Το 1 δεν είναι ούτε πρώτος ούτε σύνθετος.
\end{orismos}

Οι πρώτοι μέχρι το 50 είναι οι: \blue{2, 3, 5, 7, 11, 13, 17, 19, 23, 29, 31, 37, 41, 43, 47}. Ο 2 είναι ο μοναδικός άρτιος πρώτος, γιατί κάθε άλλος άρτιος διαιρείται και με το 2. Οι πρώτοι είναι άπειροι: το απέδειξε ο Ευκλείδης πριν από 2.300 χρόνια.

\tmimasmall{Το κόσκινο του Ερατοσθένη}

Ο Ερατοσθένης, που έζησε στην Αλεξάνδρεια τον 3ο αιώνα π.Χ., βρήκε μια μέθοδο για να «κοσκινίζουμε» τους πρώτους. Γράφεις τους αριθμούς, σβήνεις το 1, μετά κρατάς το 2 και σβήνεις όλα τα πολλαπλάσιά του, μετά κρατάς το επόμενο που δεν έσβησες, το 3, και σβήνεις τα πολλαπλάσιά του, και ούτω καθεξής. Αυτά που μένουν είναι οι πρώτοι.

\begin{center}
\renewcommand{\arraystretch}{1.25}
\setlength{\tabcolsep}{5pt}
\newcommand{\pr}[1]{\textcolor{gblue}{\bfseries #1}}
\begin{tabular}{cccccccccc}
1 & \pr{2} & \pr{3} & 4 & \pr{5} & 6 & \pr{7} & 8 & 9 & 10\\
\pr{11} & 12 & \pr{13} & 14 & 15 & 16 & \pr{17} & 18 & \pr{19} & 20\\
21 & 22 & \pr{23} & 24 & 25 & 26 & 27 & 28 & \pr{29} & 30\\
\pr{31} & 32 & 33 & 34 & 35 & 36 & \pr{37} & 38 & 39 & 40\\
\pr{41} & 42 & \pr{43} & 44 & 45 & 46 & \pr{47} & 48 & 49 & 50\\
\end{tabular}
\end{center}

\tmimasmall{Ανάλυση σε γινόμενο πρώτων παραγόντων}

Κάθε σύνθετος αριθμός γράφεται ως γινόμενο πρώτων, και \textbf{με έναν μόνο τρόπο} (αν αγνοήσουμε τη σειρά των παραγόντων). Γι' αυτό οι πρώτοι είναι τα δομικά υλικά: κάθε αριθμός έχει το δικό του «σύνολο υλικών».

\begin{kanonas}{ανάλυση σε πρώτους παράγοντες}
\begin{enumerate}[leftmargin=1.6em,itemsep=0.2ex,topsep=2pt]
\item Βρίσκω τον μικρότερο πρώτο που διαιρεί τον αριθμό (δοκιμάζω με τη σειρά 2, 3, 5, 7, 11, \dots με τα κριτήρια διαιρετότητας).
\item Διαιρώ και γράφω τον πρώτο στα δεξιά και το πηλίκο από κάτω.
\item Επαναλαμβάνω με το πηλίκο, ώσπου να φτάσω στο 1.
\item Γράφω το γινόμενο των πρώτων και τους ίδιους παράγοντες τους ενώνω σε δυνάμεις.
\end{enumerate}
\end{kanonas}

Αν ο αριθμός είναι πρώτος, δεν χρειάζεται να δοκιμάσεις όλους τους πρώτους. Δοκιμάζεις με τη σειρά και σταματάς όταν το \textbf{πηλίκο γίνει μικρότερο από τον διαιρέτη}. Αν μέχρι τότε δεν βρήκες διαιρέτη, ο αριθμός είναι πρώτος.

\begin{prosoxi}
\begin{itemize}[leftmargin=1.2em,itemsep=0.3ex,topsep=2pt]
\item Το 1 \textbf{δεν είναι πρώτος}. Αν ήταν, το 6 θα γραφόταν και $2\cdot 3$ και $1\cdot 2\cdot 3$ και $1\cdot 1\cdot 2\cdot 3$, και η ανάλυση δεν θα ήταν μοναδική.
\item Ο 2 είναι πρώτος, και μάλιστα ο μόνος άρτιος. Όχι «όλοι οι πρώτοι είναι περιττοί».
\item Ο αριθμός που μοιάζει πρώτος δεν είναι πάντα. Το 91 φαίνεται πρώτος, αλλά $91=7\cdot 13$.
\end{itemize}
\end{prosoxi}

\tmima{Λυμένα παραδείγματα}

\paradeigma{Πρώτος ή σύνθετος;}
\tixero Οι αριθμοί 91 και 97.\par
\tithelo Να αποφασίσω αν είναι πρώτοι.\par
\possindeo Για το 91: δεν διαιρείται με το 2 (τελειώνει σε 1), ούτε με το 3 ($9+1=10$), ούτε με το 5. Με το 7: $7\cdot 13=91$. Βρήκα διαιρέτη, άρα είναι σύνθετος.\par
Για το 97: ούτε με 2, ούτε με 3 ($9+7=16$), ούτε με 5. Με το 7: $97=7\cdot 13+6$, δεν διαιρείται, και το πηλίκο 13 είναι ακόμα μεγαλύτερο από το 7. Με το 11: $97=11\cdot 8+9$, και το πηλίκο 8 είναι μικρότερο από το 11. Σταματάω.\par
\apantisi Το 91 είναι σύνθετος ($7\cdot 13$), το 97 είναι πρώτος.\par
\elegxos Οι πρώτοι κάτω από το 100 είναι 25 και ο τελευταίος από αυτούς είναι το 97.

\paradeigma{Ανάλυση του 1.980}
\tixero Ο αριθμός 1.980.\par
\tithelo Να τον αναλύσω σε γινόμενο πρώτων παραγόντων.\par
\possindeo Διαιρώ διαδοχικά: $1.980:2=990$, $990:2=495$. Το 495 είναι περιττός, άρα δοκιμάζω με το 3: $495:3=165$, $165:3=55$. Μετά με το 5: $55:5=11$. Τέλος $11:11=1$.\par
Οι παράγοντες είναι δύο δυάρια, δύο τριάρια, ένα πεντάρι και ένα ενδεκάρι.\par
\apantisi $1.980=2^2\cdot 3^2\cdot 5\cdot 11$.\par
\elegxos $4\cdot 9=36$, $36\cdot 5=180$, $180\cdot 11=1.980$. Σωστά.

\paradeigma{Το μυστικό της Πολυξένης}
\tixero Το 143 είναι το γινόμενο δύο πρώτων.\par
\tithelo Ποιοι είναι αυτοί οι πρώτοι.\par
\possindeo Δοκιμάζω τους πρώτους με τη σειρά. Το 2, το 3, το 5 και το 7 δεν διαιρούν το 143. Με το 11: $143:11=13$, και το 13 είναι πρώτος.\par
\apantisi $143=11\cdot 13$.\par
\elegxos $11\cdot 13=143$. Και οι δύο παράγοντες είναι πρώτοι, άρα αυτή είναι η ανάλυση.

\tmima{Ασκήσεις}

\zoni{Το κατάλαβα}{εφαρμογή του κανόνα}

\ex Γράψε όλους τους διαιρέτες των αριθμών 16, 20, 28, 42 και 49.

\ex Γράψε τα έξι πρώτα πολλαπλάσια (μαζί με το 0) των αριθμών 3, 8 και 15.

\ex Ποιοι από τους 1, 3, 15, 23, 27, 31, 39, 43, 49, 51, 59, 63, 77 και 83 είναι πρώτοι; Για τους σύνθετους, γράψε έναν διαιρέτη τους.

\ex Βρες όλους τους πρώτους αριθμούς ανάμεσα στο 80 και στο 110.

\ex Ανάλυσε σε γινόμενο πρώτων παραγόντων (με δυνάμεις): α) 36 \quad β) 90 \quad γ) 126 \quad δ) 200 \quad ε) 360 \quad στ) 462.

\ex Ποιοι αριθμοί έχουν ανάλυση: α) $2^4\cdot 3$ \quad β) $2\cdot 5^2\cdot 7$ \quad γ) $3^2\cdot 5\cdot 7$ \quad δ) $2^3\cdot 3\cdot 11$;

\ex Σωστό ή λάθος; α) Ο 2 είναι ο μοναδικός άρτιος πρώτος. \quad β) Το 1 είναι πρώτος. \quad γ) Το 77 είναι πρώτος. \quad δ) Το άθροισμα δύο πρώτων είναι πάντα πρώτος. \quad ε) Το γινόμενο δύο πρώτων είναι πάντα σύνθετος. \quad στ) Το 0 είναι πολλαπλάσιο κάθε αριθμού. \quad ζ) Το 57 είναι πρώτος.

\zoni{Το χρησιμοποιώ}{προβλήματα από τη ζωή}

\ex \textbf{Το μυστικό της Φανής.} Η Φανή σκέφτεται δύο πρώτους με γινόμενο 323. Ποιοι είναι; Δοκίμασε με τη σειρά τους πρώτους μέχρι να βρεις διαιρέτη.

\ex \textbf{Οι καραμέλες της Μυρτώς.} Η Ρήνη έχει 54 καραμέλες και θέλει να τις βάλει σε ίσα σακουλάκια, χωρίς να περισσέψει καμία. Πόσες καραμέλες μπορεί να βάλει σε κάθε σακουλάκι; Η Ρήνη έχει και ένα άλλο βάζο με 53 καραμέλες. Με πόσους τρόπους μπορεί να τις μοιράσει εκεί; Γιατί τόσο λίγους;

\ex \textbf{Οι ομάδες.} Ο γυμναστής θέλει να χωρίσει 72 παιδιά σε ίσες ομάδες. Θέλει τουλάχιστον 2 ομάδες και τουλάχιστον 2 παιδιά σε κάθε ομάδα. Με πόσους διαφορετικούς τρόπους γίνεται; Αν τα παιδιά ήταν 73, πόσοι τρόποι θα υπήρχαν;

\ex \textbf{Η φανέλα του Ηλία.} Ο αριθμός στη φανέλα του Ηλία είναι πρώτος αριθμός ανάμεσα στο 30 και στο 50, και το άθροισμα των ψηφίων του είναι 7. Ποιος είναι ο αριθμός;

\ex \textbf{Το κόσκινο.} Φτιάξε το κόσκινο του Ερατοσθένη για τους αριθμούς από το 1 ως το 60. Πόσοι πρώτοι υπάρχουν; Πόσους αριθμούς σβήνει το 2, και πόσους ακόμα το 3;

\zoni{Το σκέφτομαι}{μια ερώτηση που πάει πιο μακριά}

\ex Βρες τους διαιρέτες του 6 και του 28, εκτός από τον ίδιο τον αριθμό, και πρόσθεσέ τους. Τι παρατηρείς; Οι μαθηματικοί λένε αυτούς τους αριθμούς «τέλειους».

\ex Δύο πρώτοι που διαφέρουν κατά 2 λέγονται δίδυμοι, όπως οι 3 και 5. Βρες όλα τα ζευγάρια διδύμων πρώτων μέχρι το 50.

\ex Ο 72 είναι $2^3\cdot 3^2$. Γράψε όλους τους αριθμούς της μορφής $2^{\alpha}\cdot 3^{\beta}$, όπου $\alpha=0,1,2,3$ και $\beta=0,1,2$. Τι είναι αυτοί οι αριθμοί για τον 72; Πόσοι είναι;

\ex Η Ζωή ρωτάει «γιατί το 1 δεν είναι πρώτος, αφού διαιρείται μόνο με τον εαυτό του και με το 1;» Εξήγησε με την ανάλυση του 6 γιατί, αν το 1 ήταν πρώτος, δεν θα είχε κάθε αριθμός μία μοναδική ανάλυση.

\ex Η Πολυξένη ισχυρίζεται ότι $10.403=101\cdot 103$. Έλεγξέ το με πολλαπλασιασμό. Πόσο δύσκολο θα ήταν να βρεις τους δύο παράγοντες χωρίς να τους ξέρεις, δοκιμάζοντας τους πρώτους με τη σειρά;

\ypotmima{Γ}{ΜΚΔ και ΕΚΠ: μοιράζω και συναντιέμαι}

\begin{anagke}
Αυτή την εβδομάδα η κυρία Αλεξάνδρα βάζει δύο προβλήματα που μοιάζουν να έχουν αντίθετη ιδέα.

Πρώτο: το σχολείο φτιάχνει κήπο. Η Α’ τάξη έχει 72 φυτά, 48 γλάστρες και 120 σακούλες χώμα. Θέλουν να φτιάξουν ίδια κιτ για τα τμήματα, ώστε κάθε κιτ να έχει και από τα τρία είδη και να μη μείνει τίποτα. «Πόσα το πολύ ίδια κιτ γίνονται;» Ο Στέλιος λέει αμέσως «δώδεκα». Η Φανή ελέγχει: δώδεκα κιτ θέλουν 6 φυτά, 4 γλάστρες και 10 σακούλες το καθένα. Βγαίνει. «Μα δεν είπαμε πόσα το πολύ;» λέει η Ζωή.

Δεύτερο: η θεία Μαρία και ο Ηλίας κάθονται το βράδυ στην ακτή και βλέπουν δύο φάρους. Ο ένας αναβοσβήνει κάθε 12 δευτερόλεπτα, ο άλλος κάθε 18. Μόλις άναψαν και οι δύο μαζί. «Σε πόσο ξανά θα ανάψουν μαζί;» ρωτάει η θεία.

Στο πρώτο θέλουμε το μεγαλύτερο κοινό κομμάτι: μοιράζουμε. Στο δεύτερο θέλουμε την πρώτη φορά που συμπίπτουν: συναντιόμαστε. Και τα δύο λύνονται με τους ίδιους αριθμούς, διαιρέτες και πολλαπλάσια.
\end{anagke}

\begin{lexi}{ΜΚΔ και ΕΚΠ}{μέγιστος, ελάχιστος, κοινός}
Οι δύο συντομογραφίες κρύβουν τρεις λέξεις η καθεμία. \lexiword{Κοινός} είναι αυτός που ανήκει και στους δύο. \lexiword{Μέγιστος} είναι ο μεγαλύτερος που υπάρχει (το «μέγας» στον υπερθετικό), και \lexiword{ελάχιστος} ο μικρότερος που υπάρχει (το «ελαχύς»).

Το \lexiword{ΜΚΔ} είναι ο Μέγιστος Κοινός Διαιρέτης: ο μεγαλύτερος αριθμός που διαιρεί και τους δύο. Το \lexiword{ΕΚΠ} είναι το Ελάχιστο Κοινό Πολλαπλάσιο: το μικρότερο (μη μηδενικό) πολλαπλάσιο που είναι και των δύο.
\end{lexi}

\begin{metafrasi}
Ο μεγαλύτερος αριθμός που διαιρεί και το 12 και το 18. & $\MKD(12,\,18)=6$\\
Ο μικρότερος αριθμός (όχι 0) που είναι πολλαπλάσιο και του 12 και του 18. & $\EKP(12,\,18)=36$\\
Οι $\alpha$ και $\beta$ δεν έχουν κοινό διαιρέτη εκτός από το 1. & $\MKD(\alpha,\beta)=1$ (πρώτοι μεταξύ τους)\\
Ο $\alpha$ διαιρεί τον $\beta$. & $\MKD(\alpha,\beta)=\alpha$, $\EKP(\alpha,\beta)=\beta$\\
Το γινόμενο των δύο αριθμών. & $\MKD\cdot\EKP=\alpha\cdot\beta$\\
\end{metafrasi}

\tmima{Η θεωρία}

\begin{orismos}{ΜΚΔ}
Ο \textbf{Μέγιστος Κοινός Διαιρέτης} δύο ή περισσότερων φυσικών είναι ο μεγαλύτερος από τους κοινούς τους διαιρέτες. Γράφουμε $\MKD(\alpha,\beta)$.
\end{orismos}

\begin{orismos}{ΕΚΠ}
Το \textbf{Ελάχιστο Κοινό Πολλαπλάσιο} δύο ή περισσότερων φυσικών είναι το μικρότερο από τα κοινά τους πολλαπλάσια, εκτός από το 0 (που είναι πολλαπλάσιο όλων). Γράφουμε $\EKP(\alpha,\beta)$.
\end{orismos}

\tmimasmall{Δύο μέθοδοι}

\textbf{Με λίστες.} Γράφω τους διαιρέτες (για ΜΚΔ) ή τα πρώτα πολλαπλάσια (για ΕΚΠ) κάθε αριθμού και διαλέγω τον κοινό που ψάχνω. Είναι εύκολο για μικρούς αριθμούς.

\begin{kanonas}{ΜΚΔ και ΕΚΠ με ανάλυση}
\begin{enumerate}[leftmargin=1.6em,itemsep=0.2ex,topsep=2pt]
\item Αναλύω κάθε αριθμό σε γινόμενο πρώτων παραγόντων (με δυνάμεις).
\item \textbf{ΜΚΔ}: παίρνω μόνο τους \emph{κοινούς} πρώτους παράγοντες, καθένα με τον \emph{μικρότερο} εκθέτη που έχει.
\item \textbf{ΕΚΠ}: παίρνω \emph{όλους} τους πρώτους παράγοντες (κοινούς και μη), καθένα με τον \emph{μεγαλύτερο} εκθέτη που έχει.
\item Πολλαπλασιάζω τους παράγοντες που πήρα.
\end{enumerate}
\end{kanonas}

Η ιδέα: το ΜΚΔ πρέπει να χωράει και στους δύο αριθμούς, γι' αυτό παίρνει όσα έχουν και οι δύο, το λιγότερο. Το ΕΚΠ πρέπει να χωράει και τους δύο, γι' αυτό παίρνει όλα όσα χρειάζεται ο καθένας, το περισσότερο.

\tmimasmall{Πρώτοι μεταξύ τους και μια σχέση}

Δύο αριθμοί λέγονται \textbf{πρώτοι μεταξύ τους} αν ο μοναδικός κοινός διαιρέτης τους είναι το 1, δηλαδή αν $\MKD(\alpha,\beta)=1$. Δεν χρειάζεται να είναι οι ίδιοι πρώτοι: οι 8 και 15 είναι σύνθετοι αλλά πρώτοι μεταξύ τους. Τότε το ΕΚΠ τους είναι το γινόμενό τους: $\EKP(8,\,15)=120$.

Για δύο φυσικούς ισχύει πάντα
$$\MKD(\alpha,\beta)\cdot\EKP(\alpha,\beta)=\alpha\cdot\beta.$$
Για τα 12 και 18: $6\cdot 36=216=12\cdot 18$.

\begin{prosoxi}
\begin{itemize}[leftmargin=1.2em,itemsep=0.3ex,topsep=2pt]
\item Το ΜΚΔ είναι πάντα μικρότερο ή ίσο από τον μικρότερο αριθμό. Το ΕΚΠ είναι πάντα μεγαλύτερο ή ίσο από τον μεγαλύτερο. Αν βγάλεις ΜΚΔ μεγαλύτερο από έναν από τους αριθμούς, κάτι πήγε στραβά.
\item Στο ΜΚΔ προσοχή: παίρνω μόνο τους \emph{κοινούς} παράγοντες. Στο ΕΚΠ παίρνω \emph{όλους}.
\item Όταν ένας αριθμός διαιρεί τον άλλον, τα πράγματα είναι εύκολα: $\MKD(7,42)=7$ και $\EKP(7,42)=42$.
\end{itemize}
\end{prosoxi}

\tmima{Λυμένα παραδείγματα}

\paradeigma{Οι φάροι}
\tixero Ο ένας φάρος αναβοσβήνει κάθε 12 δευτερόλεπτα, ο άλλος κάθε 18. Άναψαν μαζί τώρα.\par
\tithelo Σε πόσα δευτερόλεπτα θα ανάψουν ξανά μαζί.\par
\possindeo Ο πρώτος ανάβει στα δευτερόλεπτα $12, 24, 36, 48, \dots$ και ο δεύτερος στα $18, 36, 54, \dots$ Η πρώτη φορά που συμπίπτουν είναι το 36. Δηλαδή ψάχνω το $\EKP(12,\,18)$.\par
Με ανάλυση: $12=2^2\cdot 3$ και $18=2\cdot 3^2$. Παίρνω όλους τους παράγοντες με τον μεγαλύτερο εκθέτη: $2^2\cdot 3^2=36$.\par
\apantisi Σε 36 δευτερόλεπτα.\par
\elegxos Το 36 είναι πολλαπλάσιο και του 12 ($12\cdot 3$) και του 18 ($18\cdot 2$), και κανένας μικρότερος αριθμός δεν είναι (δες τις λίστες).

\paradeigma{Τα κιτ του κήπου}
\tixero 72 φυτά, 48 γλάστρες και 120 σακούλες χώμα. Τα κιτ πρέπει να είναι ίδια και να μη μείνει τίποτα.\par
\tithelo Πόσα το πολύ κιτ γίνονται και τι έχει το καθένα.\par
\possindeo Ο αριθμός των κιτ πρέπει να διαιρεί και το 72 και το 48 και το 120. Ψάχνω το μεγαλύτερο κοινό διαιρέτη. $72=2^3\cdot 3^2$, $48=2^4\cdot 3$, $120=2^3\cdot 3\cdot 5$. Κοινοί παράγοντες: το 2 (μικρότερος εκθέτης 3) και το 3 (μικρότερος εκθέτης 1). Το 5 δεν είναι κοινό. Άρα $\MKD=2^3\cdot 3=24$.\par
Κάθε κιτ έχει: $72:24=3$ φυτά, $48:24=2$ γλάστρες, $120:24=5$ σακούλες.\par
\apantisi 24 κιτ, με 3 φυτά, 2 γλάστρες και 5 σακούλες το καθένα.\par
\elegxos Τα 12 κιτ του Στέλιου βγαίνουν, αλλά το 24 είναι κοινός διαιρέτης και μεγαλύτερος από το 12. Αν δοκιμάσουμε μεγαλύτερο αριθμό, π.χ. 48 κιτ, το 72 δεν διαιρείται με το 48 ($72:48$ δεν είναι ακέραιο).

\paradeigma{ΜΚΔ και ΕΚΠ τριών αριθμών}
\tixero Οι αριθμοί 1.800, 2.100 και 3.000.\par
\tithelo Το ΜΚΔ και το ΕΚΠ τους.\par
\possindeo Αναλύω: $1.800=2^3\cdot 3^2\cdot 5^2$, $2.100=2^2\cdot 3\cdot 5^2\cdot 7$ και $3.000=2^3\cdot 3\cdot 5^3$.\par
ΜΚΔ: οι παράγοντες 2, 3 και 5 υπάρχουν και στους τρεις (το 7 όχι). Παίρνω τον μικρότερο εκθέτη του καθενός: $2^2\cdot 3\cdot 5^2=300$.\par
ΕΚΠ: παίρνω όλους τους παράγοντες με τον μεγαλύτερο εκθέτη: $2^3\cdot 3^2\cdot 5^3\cdot 7=8\cdot 9\cdot 125\cdot 7=63.000$.\par
\apantisi $\MKD=300$ και $\EKP=63.000$.\par
\elegxos Το 300 διαιρεί τον καθένα: $1.800:300=6$, $2.100:300=7$, $3.000:300=10$. Και το 63.000 είναι πολλαπλάσιο του καθενός: $63.000:1.800=35$, $63.000:2.100=30$, $63.000:3.000=21$.

\tmima{Ασκήσεις}

\zoni{Το κατάλαβα}{εφαρμογή του κανόνα}

\ex Βρες το ΜΚΔ γράφοντας τους διαιρέτες: α) 14 και 35 \quad β) 24 και 36 \quad γ) 16 και 45 \quad δ) 28 και 42.

\ex Βρες το ΕΚΠ γράφοντας τα πρώτα πολλαπλάσια: α) 6 και 8 \quad β) 9 και 12 \quad γ) 10 και 15 \quad δ) 4 και 9.

\ex Βρες το ΜΚΔ και το ΕΚΠ με ανάλυση σε πρώτους παράγοντες: α) 30 και 45 \quad β) 72 και 108 \quad γ) 84 και 126 \quad δ) 150 και 225.

\ex Βρες το ΜΚΔ και το ΕΚΠ των τριών αριθμών 20, 30 και 50.

\ex Ποια από τα ζευγάρια είναι πρώτοι μεταξύ τους; α) 9 και 16 \quad β) 15 και 28 \quad γ) 18 και 27 \quad δ) 21 και 22 \quad ε) 35 και 49 \quad στ) 1 και 50.

\ex Βρες το ΜΚΔ και το ΕΚΠ των: α) 6 και 48 \quad β) 13 και 52 \quad γ) 20 και 100. Τι παρατηρείς όταν ο ένας αριθμός διαιρεί τον άλλον;

\ex Σωστό ή λάθος; α) $\MKD(8,\,12)=24$. \quad β) $\EKP(6,\,9)=54$. \quad γ) Το ΜΚΔ δύο διαφορετικών πρώτων είναι 1. \quad δ) Το ΕΚΠ δύο αριθμών δεν είναι ποτέ μικρότερο από τον μεγαλύτερο. \quad ε) Το ΜΚΔ δεν είναι ποτέ μεγαλύτερο από τον μικρότερο. \quad στ) Αν δύο αριθμοί είναι πρώτοι μεταξύ τους, το ΕΚΠ τους είναι το γινόμενό τους.

\zoni{Το χρησιμοποιώ}{προβλήματα από τη ζωή}

\ex \textbf{Οι φάροι.} Δύο άλλοι φάροι αναβοσβήνουν ο ένας κάθε 10 δευτερόλεπτα και ο άλλος κάθε 25. Άναψαν μαζί τώρα. Πότε θα ξανανάψουν μαζί; Πόσες φορές θα συμπέσουν στα επόμενα 3 λεπτά;

\ex \textbf{Τα καλάθια της γιαγιάς.} Η γιαγιά Μυρτώ έχει 90 μήλα και 150 πορτοκάλια και θέλει να φτιάξει ίδια καλάθια, με και από τα δύο φρούτα και χωρίς να περισσέψει τίποτα. Πόσα το πολύ καλάθια μπορεί να φτιάξει; Τι έχει το καθένα;

\ex \textbf{Τα πλακάκια της θείας.} Η αυλή της θείας Μαρίας είναι ορθογώνια, 240 εκ. επί 180 εκ. Θέλει να την καλύψει με τετράγωνα πλακάκια ίδιου μεγέθους, χωρίς να κόψει κανένα. Ποια είναι η μεγαλύτερη δυνατή πλευρά του πλακιδίου; Πόσα πλακάκια θα χρειαστεί;

\ex \textbf{Τα λεωφορεία.} Δύο γραμμές λεωφορείων ξεκινούν μαζί από τη στάση στις 7:00. Η μία περνάει κάθε 12 λεπτά και η άλλη κάθε 16. Πότε θα ξαναπεράσουν μαζί; Πόσες φορές θα βρεθούν μαζί από τις 7:00 ως τις 10:00 (μετρώντας και την πρώτη φορά στις 7:00);

\ex \textbf{Το πότισμα.} Η Ρήνη ποτίζει τις γλάστρες κάθε 4 μέρες, ο Ηλίας κάθε 6 και η Ζωή κάθε 9. Σήμερα τις πότισαν και οι τρεις. Σε πόσες μέρες θα τις ποτίσουν ξανά και οι τρεις την ίδια μέρα;

\zoni{Το σκέφτομαι}{μια ερώτηση που πάει πιο μακριά}

\ex Για τα ζευγάρια 10 και 25, 16 και 24, 21 και 35, 18 και 30 βρες το ΜΚΔ, το ΕΚΠ και το γινόμενο των δύο αριθμών. Σύγκρινε το γινόμενο ΜΚΔ·ΕΚΠ με το γινόμενο των αριθμών. Τι ανακαλύπτεις;

\ex Η Ζωή παρατηρεί ότι $90=36\cdot 2+18$ και ότι $\MKD(90,\,36)=18=\MKD(36,\,18)$. Δοκίμασε το ίδιο με τα 180 και 84: γράψε την Ευκλείδεια διαίρεση και σύγκρινε το $\MKD(180,\,84)$ με το $\MKD(84,\,\upsilon)$. Αν συνεχίσεις μέχρι να βρεις υπόλοιπο 0, τι δείχνει ο τελευταίος διαιρέτης; (Αυτή είναι η μέθοδος του Ευκλείδη.)

\ex Γιατί, όταν δύο αριθμοί είναι πρώτοι μεταξύ τους, το ΕΚΠ είναι το γινόμενό τους; Εξήγησε με τα 9 και 20, χρησιμοποιώντας τις αναλύσεις τους.

\ex Βρες δύο αριθμούς που να έχουν ΜΚΔ 4 και ΕΚΠ 48. Υπάρχει και δεύτερο ζευγάρι; (Βοήθεια: το γινόμενο των δύο πρέπει να είναι $4\cdot 48$.)

\ex Μπορεί το ΜΚΔ δύο \emph{διαφορετικών} αριθμών να είναι ίσο με το ΕΚΠ τους; Εξήγησε.

\zoni{Για το σπίτι}{και λίγα ακόμη}

\ex \textbf{Ο φούρνος του χωριού.} Η γιαγιά Μυρτώ ψήνει ψωμί κάθε 5 ημέρες, ο φούρνος του χωριού καθαρίζεται κάθε 6 και το φορτηγό φέρνει αλεύρι κάθε 8. Σήμερα συνέπεσαν και τα τρία. Σε πόσες ημέρες θα συμπέσουν ξανά;

\ex \textbf{Τα κουλουράκια της Μυρτώς.} Όταν η Ρήνη βάζει τα κουλουράκια της σε σακουλάκια των 4, των 6 ή των 9, δεν περισσεύει κανένα. Πόσα κουλουράκια έχει το λιγότερο; Αν ξέρουμε ότι έχει ανάμεσα στα 150 και στα 250, πόσα μπορεί να είναι;

\ex \textbf{Μουσική και προπόνηση.} Η Αγγελική έχει μάθημα μουσικής κάθε 6 ημέρες και ο Νίκος προπόνηση κάθε 8. Μια μέρα συνέπεσαν. Σε πόσες ημέρες θα συμπέσουν ξανά; Πόσες φορές πήγε ο καθένας ξεχωριστά στο διάστημα ανάμεσα στις δύο αυτές ημέρες, χωρίς να μετρήσεις τις ημέρες που συμπίπτουν;

\ex Δύο αριθμοί έχουν ΜΚΔ το 18. Εξήγησε γιατί έχουν και άλλους κοινούς διαιρέτες εκτός από τη μονάδα.

\ex Το τριπλάσιο ενός πρώτου αριθμού είναι πρώτος ή σύνθετος; Εξήγησε.

\ex Βρες όλους τους διαιρέτες των αριθμών: α) 40 \quad β) 56 \quad γ) 66 \quad δ) 91 \quad ε) 150 \quad στ) 221.

\ex Ανάλυσε σε γινόμενο πρώτων παραγόντων τους: α) 66 \quad β) 275 \quad γ) 840 \quad δ) 1.001.

\ex Αν ένας αριθμός διαιρεί δύο αριθμούς, διαιρεί και το άθροισμα και τη διαφορά τους. Δοκίμασέ το με το 6 και τους 42 και 18. Μετά εξήγησε γιατί ο αριθμός $5\alpha+15$ διαιρείται με το 5 για κάθε φυσικό $\alpha$.

\ex Η Πολυξένη παρατηρεί ότι $37\cdot 101=3.737$. Δοκίμασε τα $24\cdot 101$ και $58\cdot 101$. Τι παρατηρείς στο αποτέλεσμα; Εξήγησε με την επιμεριστική ιδιότητα γιατί γίνεται αυτό.

\epanalipsi{1}

\tmima{Λεξικό του κεφαλαίου}

\renewcommand{\arraystretch}{1.5}
\begin{tabular*}{\linewidth}{@{\extracolsep{\fill}} >{\raggedright\arraybackslash}p{0.26\linewidth} >{\raggedright\arraybackslash}p{0.64\linewidth}}
\toprule
\textbf{Λέξη ή τύπος} & \textbf{Με μία πρόταση}\\
\midrule
φυσικός & Ένας από τους $0,\,1,2,\,3,\dots$: μετράει πράγματα που δεν σπάνε.\\
στρογγυλοποίηση & Αλλάζω έναν αριθμό με κοντινό του, με συμφωνημένα βήματα, μία φορά.\\
$\alpha^{\nu}$ & Ο $\alpha$ γραμμένος $\nu$ φορές σε γινόμενο.\\
προτεραιότητα & Παρενθέσεις, δυνάμεις, $\cdot$ και $:$, και τέλος $+$ και $-$.\\
$\Delta=\delta\cdot\pi+\upsilon$ & Ευκλείδεια διαίρεση, με $\upsilon<\delta$.\\
κριτήρια & Κανόνες στα ψηφία που δείχνουν αν ένας αριθμός διαιρείται με το 2, 3, 4, 5, 9, 10, 25.\\
πρώτος & Ο αριθμός μεγαλύτερος από το 1 με διαιρέτες μόνο τον 1 και τον εαυτό του.\\
ΜΚΔ & Ο μεγαλύτερος κοινός διαιρέτης: μοιράζω.\\
ΕΚΠ & Το μικρότερο (μη μηδενικό) κοινό πολλαπλάσιο: συναντιέμαι.\\
\bottomrule
\end{tabular*}

\tmima{Ασκήσεις επανάληψης}

\ex Γράψε με ψηφία τον αριθμό «τριάντα επτά εκατομμύρια, δύο χιλιάδες, πέντε». Μετά στρογγυλοποίησέ τον στα εκατομμύρια.

\ex Βάλε σε αύξουσα σειρά τους 5.012, 5.102, 5.021, 4.999 και 10.000.

\ex Στρογγυλοποίησε το 86.549 στις δεκάδες, στις εκατοντάδες και στις χιλιάδες.

\ex Υπολόγισε με τις ιδιότητες των πράξεων: α) $25\cdot 17\cdot 4$ \quad β) $99\cdot 8$ \quad γ) $36\cdot 15+36\cdot 85$ \quad δ) $47+68+53+32$.

\ex Βρες την τιμή: α) $3\cdot(8-2)^2-10$ \quad β) $2^5-5^2$ \quad γ) $(2+3)\cdot 4-2^3$ \quad δ) $100-4\cdot 5^2$.

\ex Κάνε την Ευκλείδεια διαίρεση $1.000:48$ και γράψε την ισότητα $\Delta=\delta\cdot\pi+\upsilon$. Ποιος αριθμός δίνει πηλίκο 7 και υπόλοιπο 6 όταν διαιρεθεί με το 15;

\ex Σήμερα είναι Πέμπτη. Τι μέρα θα είναι μετά από 100 ημέρες;

\ex Με ποιους από τους 2, 3, 4, 5, 9, 10 και 25 διαιρείται ο αριθμός 3.240;

\ex Ποιο ψηφίο πρέπει να μπει στη θέση του $\square$ ώστε ο $5\square 7$ να διαιρείται με το 9;

\ex Ποιοι από τους 57, 59, 63 και 67 είναι πρώτοι; Για τους σύνθετους, βρες την ανάλυσή τους.

\ex Ανάλυσε τον 360 σε γινόμενο πρώτων παραγόντων. Πόσους πρώτους διαφορετικούς παράγοντες έχει;

\ex Βρες το ΜΚΔ και το ΕΚΠ των αριθμών 36 και 60.

\ex \textbf{Οι κάρτες του Στέλιου.} Ο Στέλιος έχει 36 κάρτες με ζώα και 60 κάρτες με αυτοκίνητα. Θέλει να φτιάξει όσο το δυνατόν περισσότερες ίδιες στοίβες, με και από τα δύο είδη, χωρίς να περισσέψει καμία. Πόσες στοίβες θα γίνουν και τι έχει η κάθε μία;

\ex \textbf{Η πισίνα.} Ο Ηλίας πάει στην πισίνα κάθε 8 μέρες και η Ζωή κάθε 12. Σήμερα πήγαν μαζί. Πότε θα ξαναπάνε μαζί;

\ex Σωστό ή λάθος; α) Κάθε άρτιος αριθμός διαιρείται με το 4. \quad β) $\MKD(12,\,18)=6$. \quad γ) Το 91 είναι πρώτος. \quad δ) $2^4=8$. \quad ε) $17-(9-3)=17-9-3$. \quad στ) $0:5=0$.

\ex Ο Νίκος λέει ότι $6+4\cdot 5=50$ και η Ζωή ότι $6+4\cdot 5=26$. Ποιος έχει δίκιο και γιατί; Πού πρέπει να βάλεις παρένθεση ώστε να βγει 50;

\ex Πόσους διαιρέτες έχει το 100; Πόσοι από αυτούς είναι πρώτοι;

\ex \textbf{Ερωτήσεις αυτοαξιολόγησης (πράξεις και δυνάμεις).} Σωστό ή λάθος; α) $(50-20)-5=50-(20-5)$. \quad β) $5\cdot 5\cdot 5=5^3$. \quad γ) $3^4=12$. \quad δ) $\beta+\beta+\beta=3\beta$. \quad ε) $2^4+4=20$. \quad στ) $2\cdot 10^2+5\cdot 10+1=251$. \quad ζ) $30-18:6=2$. \quad η) $6\cdot 4-3+7=28$.

\ex \textbf{Ερωτήσεις αυτοαξιολόγησης (διαιρετότητα).} Σωστό ή λάθος; α) Το άθροισμα δύο περιττών αριθμών είναι πάντα άρτιος. \quad β) Αν ο $\alpha$ είναι πολλαπλάσιο του $\beta$ (με $\beta\neq 0$), τότε ο $\beta$ διαιρεί τον $\alpha$. \quad γ) Το 51 είναι πολλαπλάσιο του 3 και του 5. \quad δ) Ο 630 διαιρείται με το 9. \quad ε) Η διαίρεση $365:20$ δίνει πηλίκο 18 και υπόλοιπο 5. \quad στ) Η σχέση $250=12\cdot 20+10$ είναι Ευκλείδεια διαίρεση. \quad ζ) Ο αριθμός $4\alpha+8$ διαιρείται με το 4. \quad η) Ο 312 διαιρείται με το 4 και με το 8.

\ex \textbf{Ερωτήσεις αυτοαξιολόγησης (ΜΚΔ, ΕΚΠ, πρώτοι).} Σωστό ή λάθος; α) Ο 14 και ο 42 έχουν ΜΚΔ το 7. \quad β) $\EKP(5,\,15)=75$. \quad γ) Η γραφή $150=5\cdot 30$ είναι ανάλυση σε γινόμενο πρώτων παραγόντων. \quad δ) Δύο διαδοχικοί φυσικοί (π.χ. 14 και 15) είναι πάντα πρώτοι μεταξύ τους. \quad ε) Αν ο $\alpha$ διαιρεί τον $\beta$, τότε $\MKD(\alpha,\,\beta)=\alpha$. \quad στ) Ο 2 είναι ο μόνος άρτιος πρώτος.

\tmima{Ο χάρτης της επανάληψης}

Μετά τις λύσεις, δες σε ποιες ασκήσεις δυσκολεύτηκες. Ο πίνακας σου λέει ποιο κεφάλαιο να ξαναδιαβάσεις.

\renewcommand{\arraystretch}{1.5}
\begin{tabular*}{\linewidth}{@{\extracolsep{\fill}} >{\raggedright\arraybackslash}p{0.38\linewidth} >{\raggedright\arraybackslash}p{0.52\linewidth}}
\toprule
\textbf{Αν δυσκολεύτηκες στις} & \textbf{ξαναδιάβασε τις παραγράφους}\\
\midrule
Ε1.1, Ε1.2, Ε1.3 & 1.1 (φυσικοί, διάταξη, στρογγυλοποίηση)\\
Ε1.4 & 1.2 (ιδιότητες των πράξεων)\\
Ε1.5, Ε1.16, Ε1.18 & 1.3 (δυνάμεις, προτεραιότητα)\\
Ε1.6, Ε1.7, Ε1.19 & 1.4 (Ευκλείδεια διαίρεση)\\
Ε1.8, Ε1.9 & 1.5, μέρος Α (κριτήρια διαιρετότητας)\\
Ε1.10, Ε1.11, Ε1.17 & 1.5, μέρος Β (πρώτοι, ανάλυση)\\
Ε1.12, Ε1.13, Ε1.14 & 1.5, μέρος Γ (ΜΚΔ, ΕΚΠ)\\
Ε1.15, Ε1.18, Ε1.19, Ε1.20 & 1.1 έως 1.5\\
\bottomrule
\end{tabular*}

